\makeatletter
\def\input@path{{./}{../}{../../}{../../../}{../../../../}{preamble/}{../preamble/}{../../preamble/}{../../../preamble/}}
\makeatother

%\input{preamble_G10_T1_practice_activities_c1_c3.tex}
\input{preamble.tex}
\input{graphs/highlights.tex}
% \GTenHighlightsfalse % Uncomment to disable instructional highlights.
\input{graphs/specific_highlights.tex}

\input{graphs/G10_T1_C7_graphs.tex}
\title{10th Grade Midterm Analysis of Decisions: C7 Solutions}
\author{}
\date{}


\pagestyle{empty}

\pgfplotsset{every axis/.append style={
  xlabel={Subjective \GTenToneHL{probability}{probability} $\GTenToneHL{probability}{p}$},
  ylabel={\GTenToneHL{expected}{Expected} \GTenToneHL{expected}{Monetary} \GTenToneHL{expected}{Value}}}}
\begin{document}
	
	\maketitle
	\setcounter{tocdepth}{1}
	\setcounter{secnumdepth}{2}
	\phantomsection\label{toc}
	\vspace*{-20mm}
	\tableofcontents
	
	\section*{Evaluated Criterion}
	
	\textbf{C7}: Uses subjective probabilities and sensitivity analysis to examine how
	changes in probability assumptions affect the preferred decision.
	
	% ============================================================
	% INTRODUCTION
	% ============================================================
	\section{Introduction}\label{sec:c4-introduction}
\phantomsection\label{tone:legend}
\PracticeHighlightLegend{G10-T1-C7}

	\SectionProblemsTableOfContents{
		\SectionProblemLink{sec:c4-intro-expected-value}{Expected Value and Sensitivity Analysis} \\
		\SectionProblemLink{sec:c4-intro-2x2}{General 2x2 Problem} \\
		\SectionProblemLink{sec:c4-intro-2x3}{General 2x3 Problem} \\
		\SectionProblemLink{sec:c4-intro-3x2}{General 3x2 Problem}
	}\vspace{8mm}
	
	\subsection{Expected Value and Sensitivity Analysis}\label{sec:c4-intro-expected-value}

	
	When \GTenToneHL{probability}{probabilities} are subjective or estimated, a decision-maker can
	examine how changes in the \GTenToneHL{probability}{probability} assumptions affect the preferred
	\GTenToneHL{alternative}{alternative}. The \GTenToneHL{expected}{Expected} \GTenToneHL{expected}{Value} criterion allows each \GTenToneHL{alternative}{alternative} to be
	represented as a function of the \GTenToneHL{probability}{probability} assigned to the possible
	\GTenToneHL{state}{states} \GTenToneHL{state}{of} \GTenToneHL{state}{nature}. For a \GTenToneHL{payoff}{payoff} table with \GTenToneHL{alternative}{alternatives} in rows and \GTenToneHL{state}{states} \GTenToneHL{state}{of} \GTenToneHL{state}{nature} in columns, let
	$
	\GTenToneHL{payoff}{P_{ay}}(\GTenToneHL{alternative}{A_i},\GTenToneHL{state}{S_j})
	$
	represent the \GTenToneHL{payoff}{payoff} obtained when \GTenToneHL{alternative}{alternative} $\GTenToneHL{alternative}{A_i}$ is selected and
	\GTenToneHL{state}{state} \GTenToneHL{state}{of} \GTenToneHL{state}{nature} $\GTenToneHL{state}{S_j}$ occurs.
	
	Let $\GTenToneHL{probability}{p}$ represent the subjective \GTenToneHL{probability}{probability} that \GTenToneHL{state}{state} \GTenToneHL{state}{of} \GTenToneHL{state}{nature} $\GTenToneHL{state}{S_1}$
	occurs. Since there are two possible states,
	\begin{center}
		$\GTenToneHL{probability}{P}(\GTenToneHL{state}{S_1})=\GTenToneHL{probability}{p}$ \qquad and \qquad $\GTenToneHL{probability}{P}(\GTenToneHL{state}{S_2})=\GTenToneHL{complement}{1-p}$ \qquad with \quad $0\leq \GTenToneHL{probability}{p}\leq 1$
	\end{center}
	
	For each \GTenToneHL{alternative}{alternative}, calculate its \GTenToneHL{expected}{Expected} \GTenToneHL{expected}{Value} as a function of $\GTenToneHL{probability}{p}$:
	$$
	\GTenToneHL{expected}{EV}(\GTenToneHL{alternative}{A_i})
	=
	pP_{ay}(\GTenToneHL{alternative}{A_i},\GTenToneHL{state}{S_1})
	+
	(\GTenToneHL{complement}{1-p})\GTenToneHL{payoff}{P_{ay}}(\GTenToneHL{alternative}{A_i},\GTenToneHL{state}{S_2}).
	$$
	
	This expression shows how the expected \GTenToneHL{payoff}{payoff} changes when the
	subjective \GTenToneHL{probability}{probability} of $\GTenToneHL{state}{S_1}$ changes. The preferred \GTenToneHL{alternative}{alternative} depends on the value of $\GTenToneHL{probability}{p}$. To determine
	which \GTenToneHL{alternative}{alternative} is preferred, compare their expected-value functions.	For two \GTenToneHL{alternative}{alternatives} $\GTenToneHL{alternative}{A_1}$ and $\GTenToneHL{alternative}{A_2}$, the threshold \GTenToneHL{probability}{probability} is found
	by solving
	$$
	\GTenToneHL{expected}{EV}(\GTenToneHL{alternative}{A_1})=\GTenToneHL{expected}{EV}(\GTenToneHL{alternative}{A_2}).
	$$
	
	At this threshold, the decision-maker is indifferent between the two
	\GTenToneHL{alternative}{alternatives}. The preferred \GTenToneHL{alternative}{alternative} on each side of the threshold
	is determined by comparing the corresponding expected values. Therefore, sensitivity analysis identifies how changes in the subjective
	\GTenToneHL{probability}{probability} affect the preferred decision.
	
	\subsection{General 2x2 Problem}\label{sec:c4-intro-2x2}

	
	Consider the following \GTenToneHL{payoff}{payoff} table, where the \GTenToneHL{probability}{probability} of state
	$\GTenToneHL{state}{S_1}$ is $\GTenToneHL{probability}{p}$ and the \GTenToneHL{probability}{probability} of state $\GTenToneHL{state}{S_2}$ is $\GTenToneHL{complement}{1-p}$:
	\begin{center}
		\renewcommand{\arraystretch}{1.3}
		\begin{tabular}{lcc}
			\toprule
			Alternative
			& \begin{tabular}{c}
				$\GTenToneHL{state}{S_1}$\\
				$(\GTenToneHL{probability}{p})$
			\end{tabular}
			& \begin{tabular}{c}
				$\GTenToneHL{state}{S_2}$\\
				$(\GTenToneHL{complement}{1-p})$
			\end{tabular}
			\\
			\midrule
			$\GTenToneHL{alternative}{A_1}$
			& $\GTenToneHL{payoff}{P_{ay}}(\GTenToneHL{alternative}{A_1},\GTenToneHL{state}{S_1})$
			& $\GTenToneHL{payoff}{P_{ay}}(\GTenToneHL{alternative}{A_1},\GTenToneHL{state}{S_2})$
			\\
			$\GTenToneHL{alternative}{A_2}$
			& $\GTenToneHL{payoff}{P_{ay}}(\GTenToneHL{alternative}{A_2},\GTenToneHL{state}{S_1})$
			& $\GTenToneHL{payoff}{P_{ay}}(\GTenToneHL{alternative}{A_2},\GTenToneHL{state}{S_2})$
			\\
			\bottomrule
		\end{tabular}
	\end{center}
	
	Determine the preferred \GTenToneHL{alternative}{alternative} as a function of $\GTenToneHL{probability}{p}$, including the
	threshold \GTenToneHL{probability}{probability} at which the decision-maker is indifferent.

	\subsubsection{Solution}
	
	\textbf{Step 1: Identify the \GTenToneHL{alternative}{alternatives}, \GTenToneHL{state}{states} \GTenToneHL{state}{of} \GTenToneHL{state}{nature}, and \GTenToneHL{probability}{probabilities}}
	
	Alternatives: $\GTenToneHL{alternative}{A_1}, \GTenToneHL{alternative}{A_2}$ \qquad
	States of nature: $\GTenToneHL{state}{S_1}, \GTenToneHL{state}{S_2}$
	
	Probabilities: $\GTenToneHL{probability}{P}(\GTenToneHL{state}{S_1}) = \GTenToneHL{probability}{p}$	\qquad	$\GTenToneHL{probability}{P}(\GTenToneHL{state}{S_2}) = \GTenToneHL{complement}{1-p}$
	
	\noindent
	\setlength{\columnseprule}{0.4pt}
	
	\begin{multicols}{2}
		\vspace{0pt}
		
		\textbf{Step 2: Calculate the \GTenToneHL{expected}{expected} \GTenToneHL{expected}{value} of $\GTenToneHL{alternative}{A_1}$}
		
		For \GTenToneHL{alternative}{alternative} $\GTenToneHL{alternative}{A_1}$:
		
		\[
		\GTenToneHL{expected}{EV}(\GTenToneHL{alternative}{A_1})
		=
		pP_{ay}(\GTenToneHL{alternative}{A_1},\GTenToneHL{state}{S_1})
		+
		(\GTenToneHL{complement}{1-p})\GTenToneHL{payoff}{P_{ay}}(\GTenToneHL{alternative}{A_1},\GTenToneHL{state}{S_2}).
		\]
		
		This expression shows how the expected \GTenToneHL{payoff}{payoff} of $\GTenToneHL{alternative}{A_1}$ changes as
		$\GTenToneHL{probability}{p}$ changes.
		
		\columnbreak
		
		\textbf{Step 3: Calculate the \GTenToneHL{expected}{expected} \GTenToneHL{expected}{value} of $\GTenToneHL{alternative}{A_2}$}
		
		For \GTenToneHL{alternative}{alternative} $\GTenToneHL{alternative}{A_2}$:
		
		\[
		\GTenToneHL{expected}{EV}(\GTenToneHL{alternative}{A_2})
		=
		pP_{ay}(\GTenToneHL{alternative}{A_2},\GTenToneHL{state}{S_1})
		+
		(\GTenToneHL{complement}{1-p})\GTenToneHL{payoff}{P_{ay}}(\GTenToneHL{alternative}{A_2},\GTenToneHL{state}{S_2}).
		\]
		
		This expression shows how the expected \GTenToneHL{payoff}{payoff} of $\GTenToneHL{alternative}{A_2}$ changes as
		$\GTenToneHL{probability}{p}$ changes.
		
	\end{multicols}
	
	\setlength{\columnseprule}{0pt}
	
	\textbf{Step 4: Find the threshold \GTenToneHL{probability}{probability}}
	
	The threshold occurs when the two \GTenToneHL{alternative}{alternatives} have the same expected
	value:
	$$
	\GTenToneHL{expected}{EV}(\GTenToneHL{alternative}{A_1})=\GTenToneHL{expected}{EV}(\GTenToneHL{alternative}{A_2}).
	$$
	
	Solve this equation for $\GTenToneHL{probability}{p}$. The resulting value is the \GTenToneHL{probability}{probability} at
	which the decision-maker is indifferent between the two \GTenToneHL{alternative}{alternatives}.
	
	\textbf{Step 5: Perform the sensitivity analysis}
	
	Compare the expected values on each side of the threshold.
	
	\begin{center}
		\begin{minipage}{0.4\textwidth}
			\[
			\begin{aligned}
				\GTenToneHL{expected}{EV}(\GTenToneHL{alternative}{A_1})>\GTenToneHL{expected}{EV}(\GTenToneHL{alternative}{A_2}) &\quad\Rightarrow\quad \text{choose } \GTenToneHL{alternative}{A_1},\\
				\GTenToneHL{expected}{EV}(\GTenToneHL{alternative}{A_2})>\GTenToneHL{expected}{EV}(\GTenToneHL{alternative}{A_1}) &\quad\Rightarrow\quad \text{choose } \GTenToneHL{alternative}{A_2},\\
				\GTenToneHL{expected}{EV}(\GTenToneHL{alternative}{A_1})=\GTenToneHL{expected}{EV}(\GTenToneHL{alternative}{A_2}) &\quad\Rightarrow\quad \text{indifferent}.
			\end{aligned}
			\]
		\end{minipage}
		$ \qquad \longrightarrow \qquad$
		\begin{minipage}{0.4\textwidth}
			\textbf{Decision rule}
			\[
			\begin{cases}
				\GTenToneHL{alternative}{A_1}, & \text{if } \GTenToneHL{expected}{EV}(\GTenToneHL{alternative}{A_1})>\GTenToneHL{expected}{EV}(\GTenToneHL{alternative}{A_2}),\\
				\GTenToneHL{alternative}{A_2}, & \text{if } \GTenToneHL{expected}{EV}(\GTenToneHL{alternative}{A_2})>\GTenToneHL{expected}{EV}(\GTenToneHL{alternative}{A_1}),\\
				\text{indifferent}, & \text{if } \GTenToneHL{expected}{EV}(\GTenToneHL{alternative}{A_1})=\GTenToneHL{expected}{EV}(\GTenToneHL{alternative}{A_2}).
			\end{cases}
			\]
		\end{minipage}
	\end{center}
	
	Thus, the preferred \GTenToneHL{alternative}{alternative} can be expressed as a function of the
	subjective \GTenToneHL{probability}{probability} $\GTenToneHL{probability}{p}$.
	
	\textbf{Extended \GTenToneHL{payoff}{payoff} table}
	
	\begin{center}
		\small
		\renewcommand{\arraystretch}{1.5}
		\setlength{\tabcolsep}{4pt}
		
		\begin{tabular}{lccc}
			\toprule
			Alternative
			& $\GTenToneHL{state}{S_1}\;(\GTenToneHL{probability}{p})$
			& $\GTenToneHL{state}{S_2}\;(\GTenToneHL{complement}{1-p})$
			& \GTenToneHL{expected}{Expected} \GTenToneHL{expected}{Value}
			\\
			\midrule
			
			$\GTenToneHL{alternative}{A_1}$
			& $\GTenToneHL{payoff}{P_{ay}}(\GTenToneHL{alternative}{A_1},\GTenToneHL{state}{S_1})$
			& $\GTenToneHL{payoff}{P_{ay}}(\GTenToneHL{alternative}{A_1},\GTenToneHL{state}{S_2})$
			&
			$pP_{ay}(\GTenToneHL{alternative}{A_1},\GTenToneHL{state}{S_1})
			+
			(\GTenToneHL{complement}{1-p})\GTenToneHL{payoff}{P_{ay}}(\GTenToneHL{alternative}{A_1},\GTenToneHL{state}{S_2})$
			\\
			
			$\GTenToneHL{alternative}{A_2}$
			& $\GTenToneHL{payoff}{P_{ay}}(\GTenToneHL{alternative}{A_2},\GTenToneHL{state}{S_1})$
			& $\GTenToneHL{payoff}{P_{ay}}(\GTenToneHL{alternative}{A_2},\GTenToneHL{state}{S_2})$
			&
			$pP_{ay}(\GTenToneHL{alternative}{A_2},\GTenToneHL{state}{S_1})
			+
			(\GTenToneHL{complement}{1-p})\GTenToneHL{payoff}{P_{ay}}(\GTenToneHL{alternative}{A_2},\GTenToneHL{state}{S_2})$
			\\
			
			\bottomrule
		\end{tabular}
		
	\end{center}
	
	\subsection{General 2x3 Problem}\label{sec:c4-intro-2x3}

	
	Consider the following \GTenToneHL{payoff}{payoff} table, where the \GTenToneHL{probability}{probabilities} of the states
	of nature are subjective and may vary:
	
	\begin{multicols}{2}
		
		\begin{center}
			\renewcommand{\arraystretch}{1.3}
			\begin{tabular}{lccc}
				\toprule
				Alternative
				& \begin{tabular}{c}
					$\GTenToneHL{state}{S_1}$\\
					$(\GTenToneHL{probability}{p})$
				\end{tabular}
				& \begin{tabular}{c}
					$\GTenToneHL{state}{S_2}$\\
					$(q)$
				\end{tabular}
				& \begin{tabular}{c}
					$\GTenToneHL{state}{S_3}$\\
					$(\GTenToneHL{complement}{1-p}-q)$
				\end{tabular}
				\\
				\midrule
				$\GTenToneHL{alternative}{A_1}$
				& $\GTenToneHL{payoff}{P_{ay}}(\GTenToneHL{alternative}{A_1},\GTenToneHL{state}{S_1})$
				& $\GTenToneHL{payoff}{P_{ay}}(\GTenToneHL{alternative}{A_1},\GTenToneHL{state}{S_2})$
				& $\GTenToneHL{payoff}{P_{ay}}(\GTenToneHL{alternative}{A_1},\GTenToneHL{state}{S_3})$
				\\
				$\GTenToneHL{alternative}{A_2}$
				& $\GTenToneHL{payoff}{P_{ay}}(\GTenToneHL{alternative}{A_2},\GTenToneHL{state}{S_1})$
				& $\GTenToneHL{payoff}{P_{ay}}(\GTenToneHL{alternative}{A_2},\GTenToneHL{state}{S_2})$
				& $\GTenToneHL{payoff}{P_{ay}}(\GTenToneHL{alternative}{A_2},\GTenToneHL{state}{S_3})$
				\\
				\bottomrule
			\end{tabular}
		\end{center}
		
		\columnbreak
		
		\[
		\GTenToneHL{probability}{P}(\GTenToneHL{state}{S_1})=\GTenToneHL{probability}{p},
		\qquad
		\GTenToneHL{probability}{P}(\GTenToneHL{state}{S_2})=q,
		\qquad
		\GTenToneHL{probability}{P}(\GTenToneHL{state}{S_3})=\GTenToneHL{complement}{1-p}-q
		\]
		
		\[
		\GTenToneHL{probability}{p}\geq 0,
		\qquad
		q\geq 0,
		\qquad
		\GTenToneHL{probability}{p}+q\leq 1
		\]
		
	\end{multicols}
	
	
	Determine the preferred \GTenToneHL{alternative}{alternative} as a function of the subjective
	\GTenToneHL{probability}{probabilities} $\GTenToneHL{probability}{p}$ and $q$, including any threshold values or boundaries
	at which the preferred decision changes.
	
	\subsubsection{Solution}
	
	\textbf{Step 1: Identify the \GTenToneHL{alternative}{alternatives}, \GTenToneHL{state}{states} \GTenToneHL{state}{of} \GTenToneHL{state}{nature}, and \GTenToneHL{probability}{probabilities}}
	
	Alternatives: $\GTenToneHL{alternative}{A_1}, \GTenToneHL{alternative}{A_2}$
	\qquad
	States of nature: $\GTenToneHL{state}{S_1}, \GTenToneHL{state}{S_2}, \GTenToneHL{state}{S_3}$
	
	Probabilities:
	$
	\GTenToneHL{probability}{P}(\GTenToneHL{state}{S_1})=\GTenToneHL{probability}{p},
	\qquad
	\GTenToneHL{probability}{P}(\GTenToneHL{state}{S_2})=q,
	\qquad
	\GTenToneHL{probability}{P}(\GTenToneHL{state}{S_3})=\GTenToneHL{complement}{1-p}-q.
	$
	
	\noindent
	\setlength{\columnseprule}{0.4pt}
	
	\begin{multicols}{2}
		\vspace{0pt}
		
		\textbf{Step 2: Calculate the \GTenToneHL{expected}{expected} \GTenToneHL{expected}{value} of $\GTenToneHL{alternative}{A_1}$}
		\[
		\begin{aligned}
			\GTenToneHL{expected}{EV}(\GTenToneHL{alternative}{A_1})
			&=pP_{ay}(\GTenToneHL{alternative}{A_1},\GTenToneHL{state}{S_1})
			+qP_{ay}(\GTenToneHL{alternative}{A_1},\GTenToneHL{state}{S_2})\\
			&\quad +(\GTenToneHL{complement}{1-p}-q)\GTenToneHL{payoff}{P_{ay}}(\GTenToneHL{alternative}{A_1},\GTenToneHL{state}{S_3}).
		\end{aligned}
		\]
		
		This expression shows how the expected \GTenToneHL{payoff}{payoff} of $\GTenToneHL{alternative}{A_1}$ changes
		as the subjective \GTenToneHL{probability}{probabilities} change.
		
		\columnbreak
		
		\textbf{Step 3: Calculate the \GTenToneHL{expected}{expected} \GTenToneHL{expected}{value} of $\GTenToneHL{alternative}{A_2}$}
		\[
		\begin{aligned}
			\GTenToneHL{expected}{EV}(\GTenToneHL{alternative}{A_2})
			&=pP_{ay}(\GTenToneHL{alternative}{A_2},\GTenToneHL{state}{S_1})
			+qP_{ay}(\GTenToneHL{alternative}{A_2},\GTenToneHL{state}{S_2})\\
			&\quad +(\GTenToneHL{complement}{1-p}-q)\GTenToneHL{payoff}{P_{ay}}(\GTenToneHL{alternative}{A_2},\GTenToneHL{state}{S_3}).
		\end{aligned}
		\]
		
	\end{multicols}
	
	\setlength{\columnseprule}{0pt}
	
	\textbf{Step 4: Compare the expected-value functions and perform the sensitivity analysis}
	
	Compare
	$$
	\GTenToneHL{expected}{EV}(\GTenToneHL{alternative}{A_1}),\quad
	\GTenToneHL{expected}{EV}(\GTenToneHL{alternative}{A_2}).
	$$
	
	The preferred \GTenToneHL{alternative}{alternative} depends on the values of $\GTenToneHL{probability}{p}$ and $q$. Examine
	the preferred \GTenToneHL{alternative}{alternative} for different feasible combinations of $\GTenToneHL{probability}{p}$ and
	$q$, subject to
	$$
	\GTenToneHL{probability}{p}\geq0,\qquad q\geq0,\qquad \GTenToneHL{probability}{p}+q\leq1.
	$$
	
	The boundary occurs when the expected values of the two \GTenToneHL{alternative}{alternatives}
	are equal, indicating indifference:
	$$
	\GTenToneHL{expected}{EV}(\GTenToneHL{alternative}{A_1})=\GTenToneHL{expected}{EV}(\GTenToneHL{alternative}{A_2})
	\quad\text{at a boundary}.
	$$
	
	The preferred \GTenToneHL{alternative}{alternative} is the one with the highest \GTenToneHL{expected}{expected} \GTenToneHL{expected}{value}:
	$$
	A^*=\arg\max_{i}\GTenToneHL{expected}{EV}(\GTenToneHL{alternative}{A_i})
	$$ \vspace{-10mm}

	\textbf{Extended \GTenToneHL{payoff}{payoff} table} \vspace{-3mm}
	\begin{center}
		\small
		\renewcommand{\arraystretch}{1.5}
		\setlength{\tabcolsep}{2pt}
		
		\begin{tabular}{|l|c|c|c|c|}
			\toprule
			Alternative
			& \begin{tabular}{c}$\GTenToneHL{state}{S_1}$ \\ $\GTenToneHL{probability}{p}$\end{tabular}
			& \begin{tabular}{c}$\GTenToneHL{state}{S_2}$ \\ $q$\end{tabular}
			& \begin{tabular}{c}$\GTenToneHL{state}{S_3}$ \\ $\GTenToneHL{complement}{1-p}-q$\end{tabular}
			& \GTenToneHL{expected}{Expected} \GTenToneHL{expected}{Value}
			\\
			\midrule
			
			$\GTenToneHL{alternative}{A_1}$
			& $\GTenToneHL{payoff}{P_{ay}}(\GTenToneHL{alternative}{A_1},\GTenToneHL{state}{S_1})$
			& $\GTenToneHL{payoff}{P_{ay}}(\GTenToneHL{alternative}{A_1},\GTenToneHL{state}{S_2})$
			& $\GTenToneHL{payoff}{P_{ay}}(\GTenToneHL{alternative}{A_1},\GTenToneHL{state}{S_3})$
			&
			$pP_{ay}(\GTenToneHL{alternative}{A_1},\GTenToneHL{state}{S_1})
			+
			qP_{ay}(\GTenToneHL{alternative}{A_1},\GTenToneHL{state}{S_2})
			+
			(\GTenToneHL{complement}{1-p}-q)\GTenToneHL{payoff}{P_{ay}}(\GTenToneHL{alternative}{A_1},\GTenToneHL{state}{S_3})$
			\\
			
			$\GTenToneHL{alternative}{A_2}$
			& $\GTenToneHL{payoff}{P_{ay}}(\GTenToneHL{alternative}{A_2},\GTenToneHL{state}{S_1})$
			& $\GTenToneHL{payoff}{P_{ay}}(\GTenToneHL{alternative}{A_2},\GTenToneHL{state}{S_2})$
			& $\GTenToneHL{payoff}{P_{ay}}(\GTenToneHL{alternative}{A_2},\GTenToneHL{state}{S_3})$
			&
			$pP_{ay}(\GTenToneHL{alternative}{A_2},\GTenToneHL{state}{S_1})
			+
			qP_{ay}(\GTenToneHL{alternative}{A_2},\GTenToneHL{state}{S_2})
			+
			(\GTenToneHL{complement}{1-p}-q)\GTenToneHL{payoff}{P_{ay}}(\GTenToneHL{alternative}{A_2},\GTenToneHL{state}{S_3})$
			\\
			
			\bottomrule
		\end{tabular}
	\end{center}
	
	\subsection{General 3x2 Problem}\label{sec:c4-intro-3x2}
	
	Consider the following \GTenToneHL{payoff}{payoff} table, where the \GTenToneHL{probability}{probabilities} of the states
	of nature are subjective and may vary:
	\vspace{-4mm}
	
	\begin{multicols}{2}
		
		\begin{center}
			\renewcommand{\arraystretch}{1.3}
			\begin{tabular}{lcc}
				\toprule
				Alternative
				& \begin{tabular}{c}
					$\GTenToneHL{state}{S_1}$\\
					$(\GTenToneHL{probability}{p})$
				\end{tabular}
				& \begin{tabular}{c}
					$\GTenToneHL{state}{S_2}$\\
					$(\GTenToneHL{complement}{1-p})$
				\end{tabular}
				\\
				\midrule
				$\GTenToneHL{alternative}{A_1}$
				& $\GTenToneHL{payoff}{P_{ay}}(\GTenToneHL{alternative}{A_1},\GTenToneHL{state}{S_1})$
				& $\GTenToneHL{payoff}{P_{ay}}(\GTenToneHL{alternative}{A_1},\GTenToneHL{state}{S_2})$
				\\
				$\GTenToneHL{alternative}{A_2}$
				& $\GTenToneHL{payoff}{P_{ay}}(\GTenToneHL{alternative}{A_2},\GTenToneHL{state}{S_1})$
				& $\GTenToneHL{payoff}{P_{ay}}(\GTenToneHL{alternative}{A_2},\GTenToneHL{state}{S_2})$
				\\
				$\GTenToneHL{alternative}{A_3}$
				& $\GTenToneHL{payoff}{P_{ay}}(\GTenToneHL{alternative}{A_3},\GTenToneHL{state}{S_1})$
				& $\GTenToneHL{payoff}{P_{ay}}(\GTenToneHL{alternative}{A_3},\GTenToneHL{state}{S_2})$
				\\
				\bottomrule
			\end{tabular}
		\end{center}
		
		\columnbreak
		
		\[
		\GTenToneHL{probability}{P}(\GTenToneHL{state}{S_1})=\GTenToneHL{probability}{p},
		\qquad
		\GTenToneHL{probability}{P}(\GTenToneHL{state}{S_2})=\GTenToneHL{complement}{1-p}
		\]
		
		\[
		0\leq \GTenToneHL{probability}{p}\leq 1
		\]
		
	\end{multicols}
	
	Determine the preferred \GTenToneHL{alternative}{alternative} as a function of the subjective
	\GTenToneHL{probability}{probability} $\GTenToneHL{probability}{p}$, including any threshold values at which the preferred
	decision changes.
	
	\subsubsection{Solution}
	
	\textbf{Step 1: Identify the \GTenToneHL{alternative}{alternatives}, \GTenToneHL{state}{states} \GTenToneHL{state}{of} \GTenToneHL{state}{nature}, and \GTenToneHL{probability}{probabilities}}
	
	Alternatives: $\GTenToneHL{alternative}{A_1}, \GTenToneHL{alternative}{A_2}, \GTenToneHL{alternative}{A_3}$
	\qquad
	States of nature: $\GTenToneHL{state}{S_1}, \GTenToneHL{state}{S_2}$
	
	Probabilities:
	$
	\GTenToneHL{probability}{P}(\GTenToneHL{state}{S_1})=\GTenToneHL{probability}{p},
	\qquad
	\GTenToneHL{probability}{P}(\GTenToneHL{state}{S_2})=\GTenToneHL{complement}{1-p}.
	$
	
	\noindent
	\setlength{\columnseprule}{0.4pt}
	
	\begin{multicols}{2}
		\vspace{0pt}
		
		\textbf{Step 2: Calculate the \GTenToneHL{expected}{expected} \GTenToneHL{expected}{value} of $\GTenToneHL{alternative}{A_1}$}
		\[
		\begin{aligned}
			\GTenToneHL{expected}{EV}(\GTenToneHL{alternative}{A_1})
			&=pP_{ay}(\GTenToneHL{alternative}{A_1},\GTenToneHL{state}{S_1})\\
			&\quad +(\GTenToneHL{complement}{1-p})\GTenToneHL{payoff}{P_{ay}}(\GTenToneHL{alternative}{A_1},\GTenToneHL{state}{S_2}).
		\end{aligned}
		\]
		
		This expression shows how the expected \GTenToneHL{payoff}{payoff} of $\GTenToneHL{alternative}{A_1}$ changes
		as the subjective \GTenToneHL{probability}{probability} changes.
		
		\columnbreak
		
		\textbf{Step 3: Calculate the \GTenToneHL{expected}{expected} \GTenToneHL{expected}{value} of $\GTenToneHL{alternative}{A_2}$}
		\[
		\begin{aligned}
			\GTenToneHL{expected}{EV}(\GTenToneHL{alternative}{A_2})
			&=pP_{ay}(\GTenToneHL{alternative}{A_2},\GTenToneHL{state}{S_1})\\
			&\quad +(\GTenToneHL{complement}{1-p})\GTenToneHL{payoff}{P_{ay}}(\GTenToneHL{alternative}{A_2},\GTenToneHL{state}{S_2}).
		\end{aligned}
		\]
		
	\end{multicols}
	
	\begin{multicols}{2}
		\vspace{0pt}
		
		\textbf{Step 4: Calculate the \GTenToneHL{expected}{expected} \GTenToneHL{expected}{value} of $\GTenToneHL{alternative}{A_3}$}
		\[
		\begin{aligned}
			\GTenToneHL{expected}{EV}(\GTenToneHL{alternative}{A_3})
			&=pP_{ay}(\GTenToneHL{alternative}{A_3},\GTenToneHL{state}{S_1})\\
			&\quad +(\GTenToneHL{complement}{1-p})\GTenToneHL{payoff}{P_{ay}}(\GTenToneHL{alternative}{A_3},\GTenToneHL{state}{S_2}).
		\end{aligned}
		\]
		
		\columnbreak
		
		\textbf{Step 5: Compare the expected-value functions}
		
		Compare
		$$
		\GTenToneHL{expected}{EV}(\GTenToneHL{alternative}{A_1}),\quad
		\GTenToneHL{expected}{EV}(\GTenToneHL{alternative}{A_2}),\quad
		\GTenToneHL{expected}{EV}(\GTenToneHL{alternative}{A_3}).
		$$
		
	\end{multicols}
	
	\textbf{Step 6: Perform the sensitivity analysis}
	
	The preferred \GTenToneHL{alternative}{alternative} depends on the value of $\GTenToneHL{probability}{p}$, subject to
	$$
	0\leq \GTenToneHL{probability}{p}\leq1.
	$$
	
	For any two \GTenToneHL{alternative}{alternatives}, a threshold occurs when their expected values
	are equal, indicating indifference:
	$$
	\GTenToneHL{expected}{EV}(\GTenToneHL{alternative}{A_i})=\GTenToneHL{expected}{EV}(\GTenToneHL{alternative}{A_k})
	\quad\text{at a threshold}.
	$$
	
	The preferred \GTenToneHL{alternative}{alternative} is the one with the highest \GTenToneHL{expected}{expected} \GTenToneHL{expected}{value}:
	$$
	A^*=\arg\max_i \GTenToneHL{expected}{EV}(\GTenToneHL{alternative}{A_i})
	$$ \vspace{-10mm}
	
	\textbf{Extended \GTenToneHL{payoff}{payoff} table} \vspace{-3mm}
	\begin{center}
		\small
		\renewcommand{\arraystretch}{1.5}
		\setlength{\tabcolsep}{2pt}
		
		\begin{tabular}{|l|c|c|c|}
			\toprule
			Alternative
			& \begin{tabular}{c}$\GTenToneHL{state}{S_1}$ \\ $\GTenToneHL{probability}{p}$\end{tabular}
			& \begin{tabular}{c}$\GTenToneHL{state}{S_2}$ \\ $\GTenToneHL{complement}{1-p}$\end{tabular}
			& \GTenToneHL{expected}{Expected} \GTenToneHL{expected}{Value}
			\\
			\midrule
			
			$\GTenToneHL{alternative}{A_1}$
			& $\GTenToneHL{payoff}{P_{ay}}(\GTenToneHL{alternative}{A_1},\GTenToneHL{state}{S_1})$
			& $\GTenToneHL{payoff}{P_{ay}}(\GTenToneHL{alternative}{A_1},\GTenToneHL{state}{S_2})$
			&
			$pP_{ay}(\GTenToneHL{alternative}{A_1},\GTenToneHL{state}{S_1})
			+
			(\GTenToneHL{complement}{1-p})\GTenToneHL{payoff}{P_{ay}}(\GTenToneHL{alternative}{A_1},\GTenToneHL{state}{S_2})$
			\\
			
			$\GTenToneHL{alternative}{A_2}$
			& $\GTenToneHL{payoff}{P_{ay}}(\GTenToneHL{alternative}{A_2},\GTenToneHL{state}{S_1})$
			& $\GTenToneHL{payoff}{P_{ay}}(\GTenToneHL{alternative}{A_2},\GTenToneHL{state}{S_2})$
			&
			$pP_{ay}(\GTenToneHL{alternative}{A_2},\GTenToneHL{state}{S_1})
			+
			(\GTenToneHL{complement}{1-p})\GTenToneHL{payoff}{P_{ay}}(\GTenToneHL{alternative}{A_2},\GTenToneHL{state}{S_2})$
			\\
			
			$\GTenToneHL{alternative}{A_3}$
			& $\GTenToneHL{payoff}{P_{ay}}(\GTenToneHL{alternative}{A_3},\GTenToneHL{state}{S_1})$
			& $\GTenToneHL{payoff}{P_{ay}}(\GTenToneHL{alternative}{A_3},\GTenToneHL{state}{S_2})$
			&
			$pP_{ay}(\GTenToneHL{alternative}{A_3},\GTenToneHL{state}{S_1})
			+
			(\GTenToneHL{complement}{1-p})\GTenToneHL{payoff}{P_{ay}}(\GTenToneHL{alternative}{A_3},\GTenToneHL{state}{S_2})$
			\\
			
			\bottomrule
		\end{tabular}
	\end{center}
	
	
\phantomsection\label{tone:reference}
\subsection*{Notation reference}
\GTenToneReference{C7}
\section{Contextualized 2x2 Problem - Clothing store}
	
	A local clothing store is deciding whether to launch a new advertising campaign
	before the holiday season or continue with its current marketing strategy. The
	manager is uncertain about the level of sales during the holiday season. Let
	$\GTenToneHL{probability}{p}$ represent the subjective \GTenToneHL{probability}{probability} of strong sales. Therefore, the
	\GTenToneHL{probability}{probability} of weak sales is $\GTenToneHL{complement}{1-p}$. All \GTenToneHL{payoff}{payoffs} are expressed in thousands of
	dollars.
	
	Using the \GTenToneHL{expected}{Expected} \GTenToneHL{expected}{Value} criterion, determine how the preferred decision
	changes as the subjective \GTenToneHL{probability}{probability} $\GTenToneHL{probability}{p}$ changes. Identify the threshold
	\GTenToneHL{probability}{probability} at which the manager is indifferent between the two \GTenToneHL{alternative}{alternatives}.
	
	\vspace{-2mm}
	
	\begin{center}
		\renewcommand{\arraystretch}{1.3}
		\begin{tabular}{lcc}
			\toprule
			Alternative
			& \begin{tabular}{c}
				Strong sales $\GTenToneHL{state}{S_1}$ \\
				$\GTenToneHL{probability}{P}(\GTenToneHL{state}{S_1})=\GTenToneHL{probability}{p}$
			\end{tabular}
			& \begin{tabular}{c}
				Weak sales $\GTenToneHL{state}{S_2}$ \\
				$\GTenToneHL{probability}{P}(\GTenToneHL{state}{S_2})=\GTenToneHL{complement}{1-p}$
			\end{tabular}
			\\
			\midrule
			Launch new campaign ($\GTenToneHL{alternative}{A_1}$) & 55 & 5 \\
			Keep current strategy ($\GTenToneHL{alternative}{A_2}$) & 45 & 45 \\
			\bottomrule
		\end{tabular}
	\end{center}
	
	\vspace{-5mm}
	
	\subsection{Solution}
	
	\textbf{Step 1: Identify the \GTenToneHL{alternative}{alternatives}, \GTenToneHL{state}{states} \GTenToneHL{state}{of} \GTenToneHL{state}{nature}, and \GTenToneHL{probability}{probabilities}}
	
	Alternatives:
	$\GTenToneHL{alternative}{A_1}$ = Launch new campaign,
	$\GTenToneHL{alternative}{A_2}$ = Keep current strategy
	
	States of nature:
	$\GTenToneHL{state}{S_1}$ = Strong sales,
	$\GTenToneHL{state}{S_2}$ = Weak sales
	
	Probabilities:
	$
	\GTenToneHL{probability}{P}(\GTenToneHL{state}{S_1})=\GTenToneHL{probability}{p},
	\GTenToneHL{probability}{P}(\GTenToneHL{state}{S_2})=\GTenToneHL{complement}{1-p}
	$
	with
	$
	0\leq \GTenToneHL{probability}{p}\leq1.
	$
	
	\vspace{-3mm}
	
	\noindent
	\setlength{\columnseprule}{0.4pt}
	\begin{multicols}{2}
		\vspace{0pt}
		
		\textbf{Step 2: Calculate the \GTenToneHL{expected}{expected} \GTenToneHL{expected}{value} of $\GTenToneHL{alternative}{A_1}$}
		
		For the Launch new campaign \GTenToneHL{alternative}{alternative}:
		\[
		\GTenToneHL{expected}{EV}(\GTenToneHL{alternative}{A_1})
		=
		\GTenToneHL{probability}{p}(55)
		+
		(\GTenToneHL{complement}{1-p})(5)
		\]
		\[
		=
		55\GTenToneHL{probability}{p}+5-5\GTenToneHL{probability}{p}
		=
		\GTenToneHL{probability}{50}\GTenToneHL{probability}{p}+5.
		\]
		
		Thus, the expected \GTenToneHL{payoff}{payoff} from launching the new campaign is
		$50\GTenToneHL{probability}{p}+5$ thousand dollars.
		
		\columnbreak
		
		\textbf{Step 3: Calculate the \GTenToneHL{expected}{expected} \GTenToneHL{expected}{value} of $\GTenToneHL{alternative}{A_2}$}
		
		For the Keep current strategy \GTenToneHL{alternative}{alternative}:
		\[
		\GTenToneHL{expected}{EV}(\GTenToneHL{alternative}{A_2})
		=
		\GTenToneHL{probability}{p}(45)
		+
		(\GTenToneHL{complement}{1-p})(45)
		\]
		\[
		=
		45\GTenToneHL{probability}{p}+45-45\GTenToneHL{probability}{p}
		=
		45.
		\]
		
		Thus, the expected \GTenToneHL{payoff}{payoff} from keeping the current strategy is
		$45$ thousand dollars.
		
	\end{multicols}
	
	\EMVGraph{
		{50}{5}{blue}{$\GTenToneHL{expected}{EMV}(\GTenToneHL{alternative}{A_1})=\GTenToneHL{expected}{50}\GTenToneHL{probability}{p}+5$},
		{0}{45}{red}{$\GTenToneHL{expected}{EMV}(\GTenToneHL{alternative}{A_2})=\GTenToneHL{expected}{45}$}
	}
	
	\setlength{\columnseprule}{0pt}
	
	\textbf{Step 4: Find the threshold \GTenToneHL{probability}{probability}}
	
	The threshold occurs when the two \GTenToneHL{alternative}{alternatives} have the same expected
	value:
	$$
	\GTenToneHL{expected}{EV}(\GTenToneHL{alternative}{A_1})=\GTenToneHL{expected}{EV}(\GTenToneHL{alternative}{A_2})
	\qquad \longrightarrow \qquad
	50\GTenToneHL{probability}{p}+5=45
	$$
	$$
	\qquad \longrightarrow \qquad
	50\GTenToneHL{probability}{p}=\GTenToneHL{probability}{40}
	\qquad \longrightarrow \qquad
	\GTenToneHL{probability}{p}=0.8.
	$$
	
	Thus, the manager is indifferent between the two \GTenToneHL{alternative}{alternatives} when the
	subjective \GTenToneHL{probability}{probability} of strong sales is $80\%$.
	
	\textbf{Step 5: Perform the sensitivity analysis}
	
	Compare the expected values on each side of the threshold.
	$$
	\begin{aligned}
		\GTenToneHL{probability}{p}<0.800 &\quad\Rightarrow\quad \GTenToneHL{expected}{EV}(\GTenToneHL{alternative}{A_1})<\GTenToneHL{expected}{EV}(\GTenToneHL{alternative}{A_2}) \quad\Rightarrow\quad \text{choose } \GTenToneHL{alternative}{A_2},\\
		\GTenToneHL{probability}{p}>0.800 &\quad\Rightarrow\quad \GTenToneHL{expected}{EV}(\GTenToneHL{alternative}{A_1})>\GTenToneHL{expected}{EV}(\GTenToneHL{alternative}{A_2}) \quad\Rightarrow\quad \text{choose } \GTenToneHL{alternative}{A_1},\\
		\GTenToneHL{probability}{p}=\GTenToneHL{probability}{0.800} &\quad\Rightarrow\quad \GTenToneHL{expected}{EV}(\GTenToneHL{alternative}{A_1})=\GTenToneHL{expected}{EV}(\GTenToneHL{alternative}{A_2}) \quad\Rightarrow\quad \text{indifferent}.
	\end{aligned}
	\qquad \longrightarrow \qquad 
	\boxed{
		\begin{cases}
			\GTenToneHL{alternative}{A_2}, & 0\leq \GTenToneHL{probability}{p}<0.800,\\
			\text{indifferent}, & \GTenToneHL{probability}{p}=\GTenToneHL{probability}{0.800},\\
			\GTenToneHL{alternative}{A_1}, & 0.800<\GTenToneHL{probability}{p}\leq1.
	\end{cases}}
	$$
	
	The sensitivity analysis shows that the clothing store should keep its current
	marketing strategy when the subjective \GTenToneHL{probability}{probability} of strong sales is below
	$80\%$. Only when that \GTenToneHL{probability}{probability} exceeds $80\%$ does launching the new
	advertising campaign have the higher \GTenToneHL{expected}{expected} \GTenToneHL{expected}{value}.
	
	\textbf{Extended \GTenToneHL{payoff}{payoff} table}
	
	\begin{center}
		\renewcommand{\arraystretch}{1.5}
		\setlength{\tabcolsep}{4pt}
		
		\begin{tabular}{lccc}
			\toprule
			Alternative
			& \begin{tabular}{c}Strong sales $\GTenToneHL{state}{S_1}$ \\ $\GTenToneHL{probability}{P}(\GTenToneHL{state}{S_1})=\GTenToneHL{probability}{p}$\end{tabular}
			& \begin{tabular}{c}Weak sales $\GTenToneHL{state}{S_2}$ \\ $\GTenToneHL{probability}{P}(\GTenToneHL{state}{S_2})=\GTenToneHL{complement}{1-p}$\end{tabular}
			& \GTenToneHL{expected}{Expected} \GTenToneHL{expected}{Value}
			\\
			\midrule
			
			$\GTenToneHL{alternative}{A_1}$: Launch new campaign
			& $55$
			& $5$
			& $\GTenToneHL{probability}{p}(55)+(\GTenToneHL{complement}{1-p})(5)=50\GTenToneHL{probability}{p}+5$
			\\
			
			$\GTenToneHL{alternative}{A_2}$: Keep current strategy
			& $45$
			& $45$
			& $\GTenToneHL{probability}{p}(45)+(\GTenToneHL{complement}{1-p})(45)=45$
			\\
			
			\bottomrule
		\end{tabular}
	\end{center}
	
	
	\section{2x2 - Coffee shop}
	
	A small coffee shop is considering whether to introduce a new dessert line
	or keep its current menu. The manager is uncertain about the level of
	customer traffic during the promotional weekend. Let $\GTenToneHL{probability}{p}$ represent the
	subjective \GTenToneHL{probability}{probability} of high customer traffic. Therefore, the \GTenToneHL{probability}{probability}
	of low customer traffic is $\GTenToneHL{complement}{1-p}$. All \GTenToneHL{payoff}{payoffs} are expressed in thousands of
	dollars.
	Using the \GTenToneHL{expected}{Expected} \GTenToneHL{expected}{Value} criterion, determine how the preferred decision
	changes as the subjective \GTenToneHL{probability}{probability} $\GTenToneHL{probability}{p}$ changes. Identify the threshold
	\GTenToneHL{probability}{probability} at which the manager is indifferent between the two \GTenToneHL{alternative}{alternatives}.
	\vspace{-4mm}
	\begin{center}
		\renewcommand{\arraystretch}{1.3}
		\begin{tabular}{lcc}
			\toprule
			Alternative
			& \begin{tabular}{c}High traffic $\GTenToneHL{state}{S_1}$ \\ $\GTenToneHL{probability}{P}(\GTenToneHL{state}{S_1})=\GTenToneHL{probability}{p}$\end{tabular}
			& \begin{tabular}{c}Low traffic $\GTenToneHL{state}{S_2}$ \\ $\GTenToneHL{probability}{P}(\GTenToneHL{state}{S_2})=\GTenToneHL{complement}{1-p}$\end{tabular}
			\\
			\midrule
			Launch dessert line ($\GTenToneHL{alternative}{A_1}$) & 40 & 10 \\
			Keep current menu ($\GTenToneHL{alternative}{A_2}$) & 30 & 20 \\
			\bottomrule
		\end{tabular}
	\end{center}
	\vspace{-5mm}
	
	\subsection{Solution}
	
	\textbf{Step 1: Identify the \GTenToneHL{alternative}{alternatives}, \GTenToneHL{state}{states} \GTenToneHL{state}{of} \GTenToneHL{state}{nature}, and \GTenToneHL{probability}{probabilities}}
	
	Alternatives:
	$\GTenToneHL{alternative}{A_1}$ = Launch dessert line,
	$\GTenToneHL{alternative}{A_2}$ = Keep current menu
	
	States of nature:
	$\GTenToneHL{state}{S_1}$ = High traffic,
	$\GTenToneHL{state}{S_2}$ = Low traffic
	
	Probabilities:
	$
	\GTenToneHL{probability}{P}(\GTenToneHL{state}{S_1})=\GTenToneHL{probability}{p},
	\GTenToneHL{probability}{P}(\GTenToneHL{state}{S_2})=\GTenToneHL{complement}{1-p}
	$
	with
	$
	0\leq \GTenToneHL{probability}{p}\leq1.
	$
	
	\vspace{-3mm}
	
	\noindent
	\setlength{\columnseprule}{0.4pt} 
	\begin{multicols}{2}
		\vspace{0pt}
		
		\textbf{Step 2: Calculate the \GTenToneHL{expected}{expected} \GTenToneHL{expected}{value} of $\GTenToneHL{alternative}{A_1}$}
		
		For the Launch dessert line \GTenToneHL{alternative}{alternative}:
		\[
		\GTenToneHL{expected}{EV}(\GTenToneHL{alternative}{A_1})
		=
		\GTenToneHL{probability}{p}(40)
		+
		(\GTenToneHL{complement}{1-p})(10)
		\]
		\[
		=
		40\GTenToneHL{probability}{p}+10-10\GTenToneHL{probability}{p}
		=
		\GTenToneHL{probability}{30}\GTenToneHL{probability}{p}+10.
		\]
		
		Thus, the expected \GTenToneHL{payoff}{payoff} from launching the dessert line is
		$30\GTenToneHL{probability}{p}+10$ thousand dollars.
		
		\columnbreak
		
		\textbf{Step 3: Calculate the \GTenToneHL{expected}{expected} \GTenToneHL{expected}{value} of $\GTenToneHL{alternative}{A_2}$}
		
		For the Keep current menu \GTenToneHL{alternative}{alternative}:
		\[
		\GTenToneHL{expected}{EV}(\GTenToneHL{alternative}{A_2})
		=
		\GTenToneHL{probability}{p}(30)
		+
		(\GTenToneHL{complement}{1-p})(20)
		\]
		\[
		=
		30\GTenToneHL{probability}{p}+20-20\GTenToneHL{probability}{p}
		=
		\GTenToneHL{probability}{10}\GTenToneHL{probability}{p}+20.
		\]
		
		Thus, the expected \GTenToneHL{payoff}{payoff} from keeping the current menu is
		$10\GTenToneHL{probability}{p}+20$ thousand dollars.
		
	\end{multicols}
	
	\EMVGraph{
		{30}{10}{blue}{$\GTenToneHL{expected}{EMV}(\GTenToneHL{alternative}{A_1})=\GTenToneHL{expected}{30}\GTenToneHL{probability}{p}+10$},
		{10}{20}{red}{$\GTenToneHL{expected}{EMV}(\GTenToneHL{alternative}{A_2})=\GTenToneHL{expected}{10}\GTenToneHL{probability}{p}+20$}
	}
	
	\setlength{\columnseprule}{0pt}
	
	\textbf{Step 4: Find the threshold \GTenToneHL{probability}{probability}}
	
	The threshold occurs when the two \GTenToneHL{alternative}{alternatives} have the same expected
	value:
	\[
	\GTenToneHL{expected}{EV}(\GTenToneHL{alternative}{A_1})=\GTenToneHL{expected}{EV}(\GTenToneHL{alternative}{A_2})
	\qquad \longrightarrow \qquad
	30\GTenToneHL{probability}{p}+10=10\GTenToneHL{probability}{p}+20
	\]
	\[
	\qquad \longrightarrow \qquad
	30\GTenToneHL{probability}{p}-10\GTenToneHL{probability}{p}=\GTenToneHL{probability}{20}-10
	\qquad \longrightarrow \qquad
	20\GTenToneHL{probability}{p}=\GTenToneHL{probability}{10}
	\qquad \longrightarrow \qquad
	\GTenToneHL{probability}{p}=\frac{10}{20}
	=\frac{1}{2}
	=0.5
	\]
	
	Thus, the manager is indifferent between the two \GTenToneHL{alternative}{alternatives} when the
	subjective \GTenToneHL{probability}{probability} of high customer traffic is $50\%$.
	
	\textbf{Step 5: Perform the sensitivity analysis}
	
	Compare the expected values on each side of the threshold.
	
	\[
	\begin{aligned}
		\GTenToneHL{probability}{p}<0.5 &\quad\Rightarrow\quad \GTenToneHL{expected}{EV}(\GTenToneHL{alternative}{A_1})<\GTenToneHL{expected}{EV}(\GTenToneHL{alternative}{A_2}) \quad\Rightarrow\quad \text{choose } \GTenToneHL{alternative}{A_2},\\
		\GTenToneHL{probability}{p}>0.5 &\quad\Rightarrow\quad \GTenToneHL{expected}{EV}(\GTenToneHL{alternative}{A_1})>\GTenToneHL{expected}{EV}(\GTenToneHL{alternative}{A_2}) \quad\Rightarrow\quad \text{choose } \GTenToneHL{alternative}{A_1},\\
		\GTenToneHL{probability}{p}=\GTenToneHL{probability}{0.5} &\quad\Rightarrow\quad \GTenToneHL{expected}{EV}(\GTenToneHL{alternative}{A_1})=\GTenToneHL{expected}{EV}(\GTenToneHL{alternative}{A_2}) \quad\Rightarrow\quad \text{indifferent}.
	\end{aligned}
	\qquad \longrightarrow \qquad
	\boxed{
		\begin{cases}
			\GTenToneHL{alternative}{A_2}, & 0\leq \GTenToneHL{probability}{p}<0.5,\\
			\text{indifferent}, & \GTenToneHL{probability}{p}=\GTenToneHL{probability}{0.5},\\
			\GTenToneHL{alternative}{A_1}, & 0.5<\GTenToneHL{probability}{p}\leq1.
	\end{cases}}
	\]
	
	The sensitivity analysis shows that the coffee shop should keep its current
	menu when the subjective \GTenToneHL{probability}{probability} of high traffic is below $50\%$.
	When that \GTenToneHL{probability}{probability} exceeds $50\%$, launching the dessert line has the
	higher \GTenToneHL{expected}{expected} \GTenToneHL{expected}{value}.
	
	\textbf{Extended \GTenToneHL{payoff}{payoff} table}
	
	\begin{center}
		\renewcommand{\arraystretch}{1.5}
		\setlength{\tabcolsep}{4pt}
		
		\begin{tabular}{lccc}
			\toprule
			Alternative
			& \begin{tabular}{c}High traffic $\GTenToneHL{state}{S_1}$ \\ $\GTenToneHL{probability}{P}(\GTenToneHL{state}{S_1})=\GTenToneHL{probability}{p}$\end{tabular}
			& \begin{tabular}{c}Low traffic $\GTenToneHL{state}{S_2}$ \\ $\GTenToneHL{probability}{P}(\GTenToneHL{state}{S_2})=\GTenToneHL{complement}{1-p}$\end{tabular}
			& \GTenToneHL{expected}{Expected} \GTenToneHL{expected}{Value}
			\\
			\midrule
			
			$\GTenToneHL{alternative}{A_1}$: Launch dessert line
			& $40$ & $10$
			& $\GTenToneHL{probability}{p}(40)+(\GTenToneHL{complement}{1-p})(10)=30\GTenToneHL{probability}{p}+10$
			\\
			
			$\GTenToneHL{alternative}{A_2}$: Keep current menu
			& $30$ & $20$
			& $\GTenToneHL{probability}{p}(30)+(\GTenToneHL{complement}{1-p})(20)=10\GTenToneHL{probability}{p}+20$
			\\
			
			\bottomrule
		\end{tabular}
	\end{center}
	
	
	\section{2x2 - Manufacturing company}
	
	A small manufacturing company is deciding whether to expand its production
	capacity or maintain its current production level. The manager is uncertain
	about market demand during the next quarter. Let $\GTenToneHL{probability}{p}$ represent the
	subjective \GTenToneHL{probability}{probability} of high market demand. Therefore, the \GTenToneHL{probability}{probability}
	of low market demand is $\GTenToneHL{complement}{1-p}$. All \GTenToneHL{payoff}{payoffs} are expressed in thousands of
	dollars.
	Using the \GTenToneHL{expected}{Expected} \GTenToneHL{expected}{Value} criterion, determine which \GTenToneHL{alternative}{alternative} is preferred
	for different values of the subjective \GTenToneHL{probability}{probability} $\GTenToneHL{probability}{p}$. Determine whether
	there is a threshold \GTenToneHL{probability}{probability} within the interval $0\leq \GTenToneHL{probability}{p}\leq1$ at which
	the manager is indifferent between the two \GTenToneHL{alternative}{alternatives}.
	\vspace{-4mm}
	\begin{center}
		\renewcommand{\arraystretch}{1.3}
		\begin{tabular}{lcc}
			\toprule
			Alternative
			& \begin{tabular}{c}High demand $\GTenToneHL{state}{S_1}$ \\ $\GTenToneHL{probability}{P}(\GTenToneHL{state}{S_1})=\GTenToneHL{probability}{p}$\end{tabular}
			& \begin{tabular}{c}Low demand $\GTenToneHL{state}{S_2}$ \\ $\GTenToneHL{probability}{P}(\GTenToneHL{state}{S_2})=\GTenToneHL{complement}{1-p}$\end{tabular}
			\\
			\midrule
			Expand production ($\GTenToneHL{alternative}{A_1}$) & 55 & 30 \\
			Maintain current capacity ($\GTenToneHL{alternative}{A_2}$) & 45 & 20 \\
			\bottomrule
		\end{tabular}
	\end{center}
	\vspace{-5mm}
	
	% ============================================================
	% SOLUTION 2x2
	% ============================================================
	
	\subsection{Solution}
	
	\textbf{Step 1: Identify the \GTenToneHL{alternative}{alternatives}, \GTenToneHL{state}{states} \GTenToneHL{state}{of} \GTenToneHL{state}{nature}, and \GTenToneHL{probability}{probabilities}}
	
	Alternatives:
	$\GTenToneHL{alternative}{A_1}$ = Expand production,
	$\GTenToneHL{alternative}{A_2}$ = Maintain current capacity
	
	States of nature:
	$\GTenToneHL{state}{S_1}$ = High demand,
	$\GTenToneHL{state}{S_2}$ = Low demand
	
	Probabilities:
	$
	\GTenToneHL{probability}{P}(\GTenToneHL{state}{S_1})=\GTenToneHL{probability}{p},
	\GTenToneHL{probability}{P}(\GTenToneHL{state}{S_2})=\GTenToneHL{complement}{1-p}
	$
	with
	$
	0\leq \GTenToneHL{probability}{p}\leq1.
	$
	
	\vspace{-3mm}
	
	\noindent
	\setlength{\columnseprule}{0.4pt} 
	\begin{multicols}{2}
		\vspace{0pt}
		
		\textbf{Step 2: Calculate the \GTenToneHL{expected}{expected} \GTenToneHL{expected}{value} of $\GTenToneHL{alternative}{A_1}$}
		
		For the Expand production \GTenToneHL{alternative}{alternative}:
		\[
		\GTenToneHL{expected}{EV}(\GTenToneHL{alternative}{A_1})
		=
		\GTenToneHL{probability}{p}(55)
		+
		(\GTenToneHL{complement}{1-p})(30)
		\]
		\[
		=
		55\GTenToneHL{probability}{p}+30-30\GTenToneHL{probability}{p}
		=
		\GTenToneHL{probability}{25}\GTenToneHL{probability}{p}+30.
		\]
		
		Thus, the expected \GTenToneHL{payoff}{payoff} from expanding production is
		$25\GTenToneHL{probability}{p}+30$ thousand dollars.
		
		\columnbreak
		
		\textbf{Step 3: Calculate the \GTenToneHL{expected}{expected} \GTenToneHL{expected}{value} of $\GTenToneHL{alternative}{A_2}$}
		
		For the Maintain current capacity \GTenToneHL{alternative}{alternative}:
		\[
		\GTenToneHL{expected}{EV}(\GTenToneHL{alternative}{A_2})
		=
		\GTenToneHL{probability}{p}(45)
		+
		(\GTenToneHL{complement}{1-p})(20)
		\]
		\[
		=
		45\GTenToneHL{probability}{p}+20-20\GTenToneHL{probability}{p}
		=
		\GTenToneHL{probability}{25}\GTenToneHL{probability}{p}+20.
		\]
		
		Thus, the expected \GTenToneHL{payoff}{payoff} from maintaining current capacity is
		$25\GTenToneHL{probability}{p}+20$ thousand dollars.
		
	\end{multicols}
	
	\EMVGraph{
		{25}{30}{blue}{$\GTenToneHL{expected}{EMV}(\GTenToneHL{alternative}{A_1})=\GTenToneHL{expected}{25}\GTenToneHL{probability}{p}+30$},
		{25}{20}{red}{$\GTenToneHL{expected}{EMV}(\GTenToneHL{alternative}{A_2})=\GTenToneHL{expected}{25}\GTenToneHL{probability}{p}+20$}
	}
	
	\setlength{\columnseprule}{0pt}
	
	\textbf{Step 4: Find the threshold \GTenToneHL{probability}{probability}}
	
	The threshold occurs when the two \GTenToneHL{alternative}{alternatives} have the same expected
	value:
	\[
	\GTenToneHL{expected}{EV}(\GTenToneHL{alternative}{A_1})=\GTenToneHL{expected}{EV}(\GTenToneHL{alternative}{A_2})
	\qquad \longrightarrow \qquad
	25\GTenToneHL{probability}{p}+30=25\GTenToneHL{probability}{p}+20
	\]
	\[
	\qquad \longrightarrow \qquad
	30=20.
	\]
	
	This equation has no solution. Therefore, there is no threshold \GTenToneHL{probability}{probability}
	in the interval $0\leq \GTenToneHL{probability}{p}\leq1$ at which the manager is indifferent between
	the two \GTenToneHL{alternative}{alternatives}.
	
	\textbf{Step 5: Perform the sensitivity analysis}
	
	Compare the expected values:
	\[
	\GTenToneHL{expected}{EV}(\GTenToneHL{alternative}{A_1})-\GTenToneHL{expected}{EV}(\GTenToneHL{alternative}{A_2})
	=
	(25\GTenToneHL{probability}{p}+30)-(25\GTenToneHL{probability}{p}+20)
	=
	10.
	\]
	
	Therefore,
	\[
	\begin{aligned}
		\GTenToneHL{expected}{EV}(\GTenToneHL{alternative}{A_1})&>\GTenToneHL{expected}{EV}(\GTenToneHL{alternative}{A_2})
		\qquad \text{for all }0\leq \GTenToneHL{probability}{p}\leq1.
	\end{aligned}
	\]
	
	Thus,
	\[
	\boxed{
		\GTenToneHL{alternative}{A_1}
		\quad\text{is preferred for all}\quad
		0\leq \GTenToneHL{probability}{p}\leq1.
	}
	\]
	
	The sensitivity analysis shows that expanding production has an expected
	\GTenToneHL{payoff}{payoff} that is always $10$ thousand dollars higher than maintaining the
	current capacity, regardless of the subjective \GTenToneHL{probability}{probability} of high market
	demand.
	
	\textbf{Extended \GTenToneHL{payoff}{payoff} table}
	
	\begin{center}
		\renewcommand{\arraystretch}{1.5}
		\setlength{\tabcolsep}{4pt}
		
		\begin{tabular}{lccc}
			\toprule
			Alternative
			& \begin{tabular}{c}High demand $\GTenToneHL{state}{S_1}$ \\ $\GTenToneHL{probability}{P}(\GTenToneHL{state}{S_1})=\GTenToneHL{probability}{p}$\end{tabular}
			& \begin{tabular}{c}Low demand $\GTenToneHL{state}{S_2}$ \\ $\GTenToneHL{probability}{P}(\GTenToneHL{state}{S_2})=\GTenToneHL{complement}{1-p}$\end{tabular}
			& \GTenToneHL{expected}{Expected} \GTenToneHL{expected}{Value}
			\\
			\midrule
			
			$\GTenToneHL{alternative}{A_1}$: Expand production
			& $55$ & $30$
			& $\GTenToneHL{probability}{p}(55)+(\GTenToneHL{complement}{1-p})(30)=25\GTenToneHL{probability}{p}+30$
			\\
			
			$\GTenToneHL{alternative}{A_2}$: Maintain current capacity
			& $45$ & $20$
			& $\GTenToneHL{probability}{p}(45)+(\GTenToneHL{complement}{1-p})(20)=25\GTenToneHL{probability}{p}+20$
			\\
			
			\bottomrule
		\end{tabular}
	\end{center}
	
	\section{2x2 - Retail company}
	
	A retail company is deciding whether to launch a new store in a shopping
	district or continue operating only through its existing locations. The
	manager is uncertain about customer demand during the first year. Let $\GTenToneHL{probability}{p}$
	represent the subjective \GTenToneHL{probability}{probability} of high customer demand. Therefore,
	the \GTenToneHL{probability}{probability} of low customer demand is $\GTenToneHL{complement}{1-p}$. All \GTenToneHL{payoff}{payoffs} are expressed
	in thousands of dollars.
	
	Using the \GTenToneHL{expected}{Expected} \GTenToneHL{expected}{Value} criterion, determine which \GTenToneHL{alternative}{alternative} is preferred
	for different values of the subjective \GTenToneHL{probability}{probability} $\GTenToneHL{probability}{p}$. Determine whether
	there is a threshold \GTenToneHL{probability}{probability} within the interval $0\leq \GTenToneHL{probability}{p}\leq1$ at which
	the manager is indifferent between the two \GTenToneHL{alternative}{alternatives}.
	\vspace{-4mm}
	
	\begin{center}
		\renewcommand{\arraystretch}{1.3}
		\begin{tabular}{lcc}
			\toprule
			Alternative
			& \begin{tabular}{c}High demand $\GTenToneHL{state}{S_1}$ \\ $\GTenToneHL{probability}{P}(\GTenToneHL{state}{S_1})=\GTenToneHL{probability}{p}$\end{tabular}
			& \begin{tabular}{c}Low demand $\GTenToneHL{state}{S_2}$ \\ $\GTenToneHL{probability}{P}(\GTenToneHL{state}{S_2})=\GTenToneHL{complement}{1-p}$\end{tabular}
			\\
			\midrule
			Open new store ($\GTenToneHL{alternative}{A_1}$) & 65 & 40 \\
			Continue existing locations ($\GTenToneHL{alternative}{A_2}$) & 25 & 35 \\
			\bottomrule
		\end{tabular}
	\end{center}
	\vspace{-5mm}
	
	\subsection{Solution}
	
	\textbf{Step 1: Identify the \GTenToneHL{alternative}{alternatives}, \GTenToneHL{state}{states} \GTenToneHL{state}{of} \GTenToneHL{state}{nature}, and \GTenToneHL{probability}{probabilities}}
	
	Alternatives:
	$\GTenToneHL{alternative}{A_1}$ = Open new store,
	$\GTenToneHL{alternative}{A_2}$ = Continue existing locations
	
	States of nature:
	$\GTenToneHL{state}{S_1}$ = High customer demand,
	$\GTenToneHL{state}{S_2}$ = Low customer demand
	
	Probabilities:
	$
	\GTenToneHL{probability}{P}(\GTenToneHL{state}{S_1})=\GTenToneHL{probability}{p},
	\GTenToneHL{probability}{P}(\GTenToneHL{state}{S_2})=\GTenToneHL{complement}{1-p}
	$
	with
	$
	0\leq \GTenToneHL{probability}{p}\leq1.
	$
	
	\vspace{-3mm}
	
	\noindent
	\setlength{\columnseprule}{0.4pt} 
	\begin{multicols}{2}
		\vspace{0pt}
		
		\textbf{Step 2: Calculate the \GTenToneHL{expected}{expected} \GTenToneHL{expected}{value} of $\GTenToneHL{alternative}{A_1}$}
		
		For the Open new store \GTenToneHL{alternative}{alternative}:
		\[
		\GTenToneHL{expected}{EV}(\GTenToneHL{alternative}{A_1})
		=
		\GTenToneHL{probability}{p}(65)
		+
		(\GTenToneHL{complement}{1-p})(40)
		\]
		\[
		=
		65\GTenToneHL{probability}{p}+40-40\GTenToneHL{probability}{p}
		=
		\GTenToneHL{probability}{25}\GTenToneHL{probability}{p}+40.
		\]
		
		Thus, the expected \GTenToneHL{payoff}{payoff} from opening the new store is
		$25\GTenToneHL{probability}{p}+40$ thousand dollars.
		
		\columnbreak
		
		\textbf{Step 3: Calculate the \GTenToneHL{expected}{expected} \GTenToneHL{expected}{value} of $\GTenToneHL{alternative}{A_2}$}
		
		For the Continue existing locations \GTenToneHL{alternative}{alternative}:
		\[
		\GTenToneHL{expected}{EV}(\GTenToneHL{alternative}{A_2})
		=
		\GTenToneHL{probability}{p}(25)
		+
		(\GTenToneHL{complement}{1-p})(35)
		\]
		\[
		=
		25\GTenToneHL{probability}{p}+35-35\GTenToneHL{probability}{p}
		=
		\GTenToneHL{probability}{35}-10\GTenToneHL{probability}{p}.
		\]
		
		Thus, the expected \GTenToneHL{payoff}{payoff} from continuing with the existing locations is
		$35-10\GTenToneHL{probability}{p}$ thousand dollars.
		
	\end{multicols}
	
	\EMVGraph{
		{25}{40}{blue}{$\GTenToneHL{expected}{EMV}(\GTenToneHL{alternative}{A_1})=\GTenToneHL{expected}{25}\GTenToneHL{probability}{p}+40$},
		{-10}{35}{red}{$\GTenToneHL{expected}{EMV}(\GTenToneHL{alternative}{A_2})=\GTenToneHL{expected}{35}-10\GTenToneHL{probability}{p}$}
	}
	
	\setlength{\columnseprule}{0pt}
	
	\textbf{Step 4: Find the threshold \GTenToneHL{probability}{probability}}
	
	The threshold occurs when the two \GTenToneHL{alternative}{alternatives} have the same expected
	value:
	\[
	\GTenToneHL{expected}{EV}(\GTenToneHL{alternative}{A_1})=\GTenToneHL{expected}{EV}(\GTenToneHL{alternative}{A_2})
	\qquad \longrightarrow \qquad
	25\GTenToneHL{probability}{p}+40=35-10\GTenToneHL{probability}{p}
	\]
	\[
	\qquad \longrightarrow \qquad
	35\GTenToneHL{probability}{p}=\GTenToneHL{probability}{-5}
	\qquad \longrightarrow \qquad
	\GTenToneHL{probability}{p}=-\frac{1}{7}
	\approx -0.143.
	\]
	
	The threshold \GTenToneHL{probability}{probability} is negative and therefore lies outside the
	feasible interval $0\leq \GTenToneHL{probability}{p}\leq1$. Thus, there is no threshold \GTenToneHL{probability}{probability}
	within the feasible interval at which the manager is indifferent between
	the two \GTenToneHL{alternative}{alternatives}.
	
	\textbf{Step 5: Perform the sensitivity analysis}
	
	Compare the expected values on the feasible interval.
	
	\[
	\GTenToneHL{expected}{EV}(\GTenToneHL{alternative}{A_1})-\GTenToneHL{expected}{EV}(\GTenToneHL{alternative}{A_2})
	=
	(25\GTenToneHL{probability}{p}+40)-(35-10\GTenToneHL{probability}{p})
	=
	35\GTenToneHL{probability}{p}+5.
	\]
	
	Since $35\GTenToneHL{probability}{p}+5>0$ for every $0\leq \GTenToneHL{probability}{p}\leq1$,
	\[
	\boxed{
		\GTenToneHL{alternative}{A_1}
		\quad\text{is preferred for all}\quad
		0\leq \GTenToneHL{probability}{p}\leq1.
	}
	\]
	
	The sensitivity analysis shows that opening the new store has the higher
	expected \GTenToneHL{payoff}{payoff} for every feasible subjective \GTenToneHL{probability}{probability} of high customer
	demand.
	
	\textbf{Extended \GTenToneHL{payoff}{payoff} table}
	
	\begin{center}
		\renewcommand{\arraystretch}{1.5}
		\setlength{\tabcolsep}{4pt}
		
		\begin{tabular}{lccc}
			\toprule
			Alternative
			& \begin{tabular}{c}High demand $\GTenToneHL{state}{S_1}$ \\ $\GTenToneHL{probability}{P}(\GTenToneHL{state}{S_1})=\GTenToneHL{probability}{p}$\end{tabular}
			& \begin{tabular}{c}Low demand $\GTenToneHL{state}{S_2}$ \\ $\GTenToneHL{probability}{P}(\GTenToneHL{state}{S_2})=\GTenToneHL{complement}{1-p}$\end{tabular}
			& \GTenToneHL{expected}{Expected} \GTenToneHL{expected}{Value}
			\\
			\midrule
			
			$\GTenToneHL{alternative}{A_1}$: Open new store
			& $65$ & $40$
			& $\GTenToneHL{probability}{p}(65)+(\GTenToneHL{complement}{1-p})(40)=25\GTenToneHL{probability}{p}+40$
			\\
			
			$\GTenToneHL{alternative}{A_2}$: Continue existing locations
			& $25$ & $35$
			& $\GTenToneHL{probability}{p}(25)+(\GTenToneHL{complement}{1-p})(35)=35-10\GTenToneHL{probability}{p}$
			\\
			
			\bottomrule
		\end{tabular}
	\end{center}
	
	\section{2x2 - Retail company}
	
	A retail company is deciding whether to launch a new store in a shopping
	district or continue operating only through its existing locations. The
	manager is uncertain about customer demand during the first year. Let $\GTenToneHL{probability}{p}$
	represent the subjective \GTenToneHL{probability}{probability} of high customer demand. Therefore,
	the \GTenToneHL{probability}{probability} of low customer demand is $\GTenToneHL{complement}{1-p}$. All \GTenToneHL{payoff}{payoffs} are expressed
	in thousands of dollars.
	
	Using the \GTenToneHL{expected}{Expected} \GTenToneHL{expected}{Value} criterion, determine which \GTenToneHL{alternative}{alternative} is preferred
	for different values of the subjective \GTenToneHL{probability}{probability} $\GTenToneHL{probability}{p}$. Determine whether
	there is a threshold \GTenToneHL{probability}{probability} within the interval $0\leq \GTenToneHL{probability}{p}\leq1$ at which
	the manager is indifferent between the two \GTenToneHL{alternative}{alternatives}.
	\vspace{-4mm}
	
	\begin{center}
		\renewcommand{\arraystretch}{1.3}
		\begin{tabular}{lcc}
			\toprule
			Alternative
			& \begin{tabular}{c}High demand $\GTenToneHL{state}{S_1}$ \\ $\GTenToneHL{probability}{P}(\GTenToneHL{state}{S_1})=\GTenToneHL{probability}{p}$\end{tabular}
			& \begin{tabular}{c}Low demand $\GTenToneHL{state}{S_2}$ \\ $\GTenToneHL{probability}{P}(\GTenToneHL{state}{S_2})=\GTenToneHL{complement}{1-p}$\end{tabular}
			\\
			\midrule
			Open new store ($\GTenToneHL{alternative}{A_1}$) & 40 & 65 \\
			Continue existing locations ($\GTenToneHL{alternative}{A_2}$) & 35 & 25 \\
			\bottomrule
		\end{tabular}
	\end{center}
	\vspace{-5mm}
	
	% ============================================================
	% SOLUTION 2x2
	% ============================================================
	
	\subsection{Solution}
	
	\textbf{Step 1: Identify the \GTenToneHL{alternative}{alternatives}, \GTenToneHL{state}{states} \GTenToneHL{state}{of} \GTenToneHL{state}{nature}, and \GTenToneHL{probability}{probabilities}}
	
	Alternatives:
	$\GTenToneHL{alternative}{A_1}$ = Open new store,
	$\GTenToneHL{alternative}{A_2}$ = Continue existing locations
	
	States of nature:
	$\GTenToneHL{state}{S_1}$ = High customer demand,
	$\GTenToneHL{state}{S_2}$ = Low customer demand
	
	Probabilities:
	$
	\GTenToneHL{probability}{P}(\GTenToneHL{state}{S_1})=\GTenToneHL{probability}{p},
	\GTenToneHL{probability}{P}(\GTenToneHL{state}{S_2})=\GTenToneHL{complement}{1-p}
	$
	with
	$
	0\leq \GTenToneHL{probability}{p}\leq1.
	$
	
	\vspace{-3mm}
	
	\noindent
	\setlength{\columnseprule}{0.4pt} 
	\begin{multicols}{2}
		\vspace{0pt}
		
		\textbf{Step 2: Calculate the \GTenToneHL{expected}{expected} \GTenToneHL{expected}{value} of $\GTenToneHL{alternative}{A_1}$}
		
		For the Open new store \GTenToneHL{alternative}{alternative}:
		\[
		\GTenToneHL{expected}{EV}(\GTenToneHL{alternative}{A_1})
		=
		\GTenToneHL{probability}{p}(40)
		+
		(\GTenToneHL{complement}{1-p})(65)
		\]
		\[
		=
		40\GTenToneHL{probability}{p}+65-65\GTenToneHL{probability}{p}
		=
		\GTenToneHL{probability}{65}-25\GTenToneHL{probability}{p}.
		\]
		
		Thus, the expected \GTenToneHL{payoff}{payoff} from opening the new store is
		$65-25\GTenToneHL{probability}{p}$ thousand dollars.
		
		\columnbreak
		
		\textbf{Step 3: Calculate the \GTenToneHL{expected}{expected} \GTenToneHL{expected}{value} of $\GTenToneHL{alternative}{A_2}$}
		
		For the Continue existing locations \GTenToneHL{alternative}{alternative}:
		\[
		\GTenToneHL{expected}{EV}(\GTenToneHL{alternative}{A_2})
		=
		\GTenToneHL{probability}{p}(35)
		+
		(\GTenToneHL{complement}{1-p})(25)
		\]
		\[
		=
		35\GTenToneHL{probability}{p}+25-25\GTenToneHL{probability}{p}
		=
		\GTenToneHL{probability}{10}\GTenToneHL{probability}{p}+25.
		\]
		
		Thus, the expected \GTenToneHL{payoff}{payoff} from continuing with the existing locations is
		$10\GTenToneHL{probability}{p}+25$ thousand dollars.
		
	\end{multicols}
	
	\EMVGraph{
		{-25}{65}{blue}{$\GTenToneHL{expected}{EMV}(\GTenToneHL{alternative}{A_1})=\GTenToneHL{expected}{65}-25\GTenToneHL{probability}{p}$},
		{10}{25}{red}{$\GTenToneHL{expected}{EMV}(\GTenToneHL{alternative}{A_2})=\GTenToneHL{expected}{10}\GTenToneHL{probability}{p}+25$}
	}
	
	\setlength{\columnseprule}{0pt}
	
	\textbf{Step 4: Find the threshold \GTenToneHL{probability}{probability}}
	
	The threshold occurs when the two \GTenToneHL{alternative}{alternatives} have the same expected
	value:
	\[
	\GTenToneHL{expected}{EV}(\GTenToneHL{alternative}{A_1})=\GTenToneHL{expected}{EV}(\GTenToneHL{alternative}{A_2})
	\qquad \longrightarrow \qquad
	65-25\GTenToneHL{probability}{p}=\GTenToneHL{probability}{10}\GTenToneHL{probability}{p}+25
	\]
	\[
	\qquad \longrightarrow \qquad
	40=35\GTenToneHL{probability}{p}
	\qquad \longrightarrow \qquad
	\GTenToneHL{probability}{p}=\frac{40}{35}
	=\frac{8}{7}
	\approx1.143.
	\]
	
	The threshold \GTenToneHL{probability}{probability} is greater than $1$ and therefore lies outside
	the feasible interval $0\leq \GTenToneHL{probability}{p}\leq1$. Thus, there is no threshold \GTenToneHL{probability}{probability}
	within the feasible interval at which the manager is indifferent between
	the two \GTenToneHL{alternative}{alternatives}.
	
	\textbf{Step 5: Perform the sensitivity analysis}
	
	Compare the expected values on the feasible interval.
	
	\[
	\GTenToneHL{expected}{EV}(\GTenToneHL{alternative}{A_1})-\GTenToneHL{expected}{EV}(\GTenToneHL{alternative}{A_2})
	=
	(65-25\GTenToneHL{probability}{p})-(10\GTenToneHL{probability}{p}+25)
	=
	40-35\GTenToneHL{probability}{p}.
	\]
	
	Since
	\[
	40-35\GTenToneHL{probability}{p}\geq5>0
	\qquad\text{for all }0\leq \GTenToneHL{probability}{p}\leq1,
	\]
	we have
	\[
	\boxed{
		\GTenToneHL{alternative}{A_1}
		\quad\text{is preferred for all}\quad
		0\leq \GTenToneHL{probability}{p}\leq1.
	}
	\]
	
	The sensitivity analysis shows that opening the new store has the higher
	expected \GTenToneHL{payoff}{payoff} for every feasible subjective \GTenToneHL{probability}{probability} of high customer
	demand.
	
	\textbf{Extended \GTenToneHL{payoff}{payoff} table}
	
	\begin{center}
		\renewcommand{\arraystretch}{1.5}
		\setlength{\tabcolsep}{4pt}
		
		\begin{tabular}{lccc}
			\toprule
			Alternative
			& \begin{tabular}{c}High demand $\GTenToneHL{state}{S_1}$ \\ $\GTenToneHL{probability}{P}(\GTenToneHL{state}{S_1})=\GTenToneHL{probability}{p}$\end{tabular}
			& \begin{tabular}{c}Low demand $\GTenToneHL{state}{S_2}$ \\ $\GTenToneHL{probability}{P}(\GTenToneHL{state}{S_2})=\GTenToneHL{complement}{1-p}$\end{tabular}
			& \GTenToneHL{expected}{Expected} \GTenToneHL{expected}{Value}
			\\
			\midrule
			
			$\GTenToneHL{alternative}{A_1}$: Open new store
			& $40$ & $65$
			& $\GTenToneHL{probability}{p}(40)+(\GTenToneHL{complement}{1-p})(65)=65-25\GTenToneHL{probability}{p}$
			\\
			
			$\GTenToneHL{alternative}{A_2}$: Continue existing locations
			& $35$ & $25$
			& $\GTenToneHL{probability}{p}(35)+(\GTenToneHL{complement}{1-p})(25)=10\GTenToneHL{probability}{p}+25$
			\\
			
			\bottomrule
		\end{tabular}
	\end{center}
	
	\section{2x3 - Delivery company}
	
	A delivery company is deciding whether to expand its delivery fleet or continue
	operating with its current fleet. The manager is uncertain about the level of
	demand during the next year. Demand may be high, moderate, or low.
	
	Let $\GTenToneHL{probability}{p}$ represent the subjective \GTenToneHL{probability}{probability} of high demand. The \GTenToneHL{probability}{probability}
	of moderate demand is $2\GTenToneHL{probability}{p}$, while the \GTenToneHL{probability}{probability} of low demand is $1-3\GTenToneHL{probability}{p}$.
	All \GTenToneHL{payoff}{payoffs} are expressed in thousands of dollars.
	
	Using the \GTenToneHL{expected}{Expected} \GTenToneHL{expected}{Value} criterion, determine which \GTenToneHL{alternative}{alternative} is preferred
	for different values of the subjective \GTenToneHL{probability}{probability} $\GTenToneHL{probability}{p}$. Identify the threshold
	\GTenToneHL{probability}{probability} at which the manager is indifferent between the two \GTenToneHL{alternative}{alternatives}.
	
	\vspace{-4mm}
	
	\begin{center}
		\renewcommand{\arraystretch}{1.3}
		\begin{tabular}{lccc}
			\toprule
			Alternative
			& \begin{tabular}{c}High demand $\GTenToneHL{state}{S_1}$ \\ $\GTenToneHL{probability}{P}(\GTenToneHL{state}{S_1})=\GTenToneHL{probability}{p}$\end{tabular}
			& \begin{tabular}{c}Moderate demand $\GTenToneHL{state}{S_2}$ \\ $\GTenToneHL{probability}{P}(\GTenToneHL{state}{S_2})=\GTenToneHL{probability}{2}\GTenToneHL{probability}{p}$\end{tabular}
			& \begin{tabular}{c}Low demand $\GTenToneHL{state}{S_3}$ \\ $\GTenToneHL{probability}{P}(\GTenToneHL{state}{S_3})=\GTenToneHL{probability}{1}-3\GTenToneHL{probability}{p}$\end{tabular}
			\\
			\midrule
			Expand fleet ($\GTenToneHL{alternative}{A_1}$) & 80 & 55 & 20 \\
			Keep current fleet ($\GTenToneHL{alternative}{A_2}$) & 50 & 45 & 35 \\
			\bottomrule
		\end{tabular}
	\end{center}
	
	\vspace{-5mm}
	
	% ============================================================
	% SOLUTION 2x3
	% ============================================================
	
	\subsection{Solution}
	
	\textbf{Step 1: Identify the \GTenToneHL{alternative}{alternatives}, \GTenToneHL{state}{states} \GTenToneHL{state}{of} \GTenToneHL{state}{nature}, and \GTenToneHL{probability}{probabilities}}
	
	Alternatives:
	$\GTenToneHL{alternative}{A_1}$ = Expand fleet,
	$\GTenToneHL{alternative}{A_2}$ = Keep current fleet
	
	States of nature:
	$\GTenToneHL{state}{S_1}$ = High demand,
	$\GTenToneHL{state}{S_2}$ = Moderate demand,
	$\GTenToneHL{state}{S_3}$ = Low demand
	
	Probabilities:
	\[
	\GTenToneHL{probability}{P}(\GTenToneHL{state}{S_1})=\GTenToneHL{probability}{p},
	\qquad
	\GTenToneHL{probability}{P}(\GTenToneHL{state}{S_2})=\GTenToneHL{probability}{2}\GTenToneHL{probability}{p},
	\qquad
	\GTenToneHL{probability}{P}(\GTenToneHL{state}{S_3})=\GTenToneHL{probability}{1}-3\GTenToneHL{probability}{p}.
	\]
	
	Since all \GTenToneHL{probability}{probabilities} must be nonnegative,
	\[
	1-3\GTenToneHL{probability}{p}\geq0
	\qquad\Longrightarrow\qquad
	\GTenToneHL{probability}{p}\leq\frac13.
	\]
	
	Therefore, the feasible interval is
	\[
	0\leq \GTenToneHL{probability}{p}\leq\frac13.
	\]
	
	\vspace{-3mm}
	
	\noindent
	\setlength{\columnseprule}{0.4pt}
	\begin{multicols}{2}
		\vspace{0pt}
		
		\textbf{Step 2: Calculate the \GTenToneHL{expected}{expected} \GTenToneHL{expected}{value} of $\GTenToneHL{alternative}{A_1}$}
		
		For the Expand fleet \GTenToneHL{alternative}{alternative}:
		\[
		\GTenToneHL{expected}{EV}(\GTenToneHL{alternative}{A_1})
		=
		\GTenToneHL{probability}{p}(80)+2\GTenToneHL{probability}{p}(55)+(1-3\GTenToneHL{probability}{p})(20)
		\]
		\[
		=
		80\GTenToneHL{probability}{p}+110\GTenToneHL{probability}{p}+20-60\GTenToneHL{probability}{p}
		=
		\GTenToneHL{probability}{130}\GTenToneHL{probability}{p}+20.
		\]
		
		Thus, the expected \GTenToneHL{payoff}{payoff} from expanding the fleet is
		$130\GTenToneHL{probability}{p}+20$ thousand dollars.
		
		\columnbreak
		
		\textbf{Step 3: Calculate the \GTenToneHL{expected}{expected} \GTenToneHL{expected}{value} of $\GTenToneHL{alternative}{A_2}$}
		
		For the Keep current fleet \GTenToneHL{alternative}{alternative}:
		\[
		\GTenToneHL{expected}{EV}(\GTenToneHL{alternative}{A_2})
		=
		\GTenToneHL{probability}{p}(50)+2\GTenToneHL{probability}{p}(45)+(1-3\GTenToneHL{probability}{p})(35)
		\]
		\[
		=
		50\GTenToneHL{probability}{p}+90\GTenToneHL{probability}{p}+35-105\GTenToneHL{probability}{p}
		=
		\GTenToneHL{probability}{35}\GTenToneHL{probability}{p}+35.
		\]
		
		Thus, the expected \GTenToneHL{payoff}{payoff} from keeping the current fleet is
		$35\GTenToneHL{probability}{p}+35$ thousand dollars.
		
	\end{multicols}
	
	\EMVGraph{
		{130}{20}{blue}{$\GTenToneHL{expected}{EMV}(\GTenToneHL{alternative}{A_1})=\GTenToneHL{expected}{130}\GTenToneHL{probability}{p}+20$},
		{35}{35}{red}{$\GTenToneHL{expected}{EMV}(\GTenToneHL{alternative}{A_2})=\GTenToneHL{expected}{35}\GTenToneHL{probability}{p}+35$}
	}
	
	\setlength{\columnseprule}{0pt}
	
	\textbf{Step 4: Find the threshold \GTenToneHL{probability}{probability}}
	
	The threshold occurs when the two \GTenToneHL{alternative}{alternatives} have the same \GTenToneHL{expected}{expected} \GTenToneHL{expected}{value}:
	\[
	\GTenToneHL{expected}{EV}(\GTenToneHL{alternative}{A_1})=\GTenToneHL{expected}{EV}(\GTenToneHL{alternative}{A_2})
	\qquad\longrightarrow\qquad
	130\GTenToneHL{probability}{p}+20=35\GTenToneHL{probability}{p}+35
	\]
	\[
	\qquad\longrightarrow\qquad
	95\GTenToneHL{probability}{p}=\GTenToneHL{probability}{15}
	\qquad\longrightarrow\qquad
	\GTenToneHL{probability}{p}=\frac{15}{95}
	=\frac{3}{19}
	\approx0.158.
	\]
	
	Since
	\[
	0<\frac{3}{19}<\frac13,
	\]
	the threshold lies within the feasible interval.
	
	\textbf{Step 5: Perform the sensitivity analysis}
	
	Compare the expected values:
	\[
	\GTenToneHL{expected}{EV}(\GTenToneHL{alternative}{A_1})-\GTenToneHL{expected}{EV}(\GTenToneHL{alternative}{A_2})
	=
	(130\GTenToneHL{probability}{p}+20)-(35\GTenToneHL{probability}{p}+35)
	=
	95\GTenToneHL{probability}{p}-15.
	\]
	
	Therefore,
	\[
	\boxed{
		\begin{aligned}
			\GTenToneHL{alternative}{A_2} &\text{ is preferred if } 0\leq \GTenToneHL{probability}{p}<\frac{3}{19},\\
			\GTenToneHL{alternative}{A_1} &\text{ and }\GTenToneHL{alternative}{A_2}\text{ are equivalent if }\GTenToneHL{probability}{p}=\frac{3}{19},\\
			\GTenToneHL{alternative}{A_1} &\text{ is preferred if }\frac{3}{19}<\GTenToneHL{probability}{p}\leq\frac13.
		\end{aligned}
	}
	\]
	
	The sensitivity analysis shows that keeping the current fleet is preferred
	when the subjective \GTenToneHL{probability}{probability} of high demand is relatively low. Once $\GTenToneHL{probability}{p}$
	exceeds $\frac{3}{19}$, expanding the fleet has the higher expected \GTenToneHL{payoff}{payoff}.
	
	\textbf{Extended \GTenToneHL{payoff}{payoff} table}
	
	\begin{center}
		\renewcommand{\arraystretch}{1.5}
		\setlength{\tabcolsep}{4pt}
		
		\begin{tabular}{lcccc}
			\toprule
			Alternative
			& \begin{tabular}{c}High demand $\GTenToneHL{state}{S_1}$ \\ $\GTenToneHL{probability}{P}(\GTenToneHL{state}{S_1})=\GTenToneHL{probability}{p}$\end{tabular}
			& \begin{tabular}{c}Moderate demand $\GTenToneHL{state}{S_2}$ \\ $\GTenToneHL{probability}{P}(\GTenToneHL{state}{S_2})=\GTenToneHL{probability}{2}\GTenToneHL{probability}{p}$\end{tabular}
			& \begin{tabular}{c}Low demand $\GTenToneHL{state}{S_3}$ \\ $\GTenToneHL{probability}{P}(\GTenToneHL{state}{S_3})=\GTenToneHL{probability}{1}-3\GTenToneHL{probability}{p}$\end{tabular}
			& \GTenToneHL{expected}{Expected} \GTenToneHL{expected}{Value}
			\\
			\midrule
			
			$\GTenToneHL{alternative}{A_1}$: Expand fleet
			& $80$ & $55$ & $20$
			& $\GTenToneHL{probability}{p}(80)+2\GTenToneHL{probability}{p}(55)+(1-3\GTenToneHL{probability}{p})(20)=130\GTenToneHL{probability}{p}+20$
			\\
			
			$\GTenToneHL{alternative}{A_2}$: Keep current fleet
			& $50$ & $45$ & $35$
			& $\GTenToneHL{probability}{p}(50)+2\GTenToneHL{probability}{p}(45)+(1-3\GTenToneHL{probability}{p})(35)=35\GTenToneHL{probability}{p}+35$
			\\
			
			\bottomrule
		\end{tabular}
	\end{center}
	
	\section{2x3 - Manufacturing company}
	
	A manufacturing company is deciding whether to invest in a new automated
	production system or continue using its current production system. The manager
	is uncertain about market conditions during the next year. Market conditions
	may be strong, moderate, or weak.
	
	Let $\GTenToneHL{probability}{p}$ represent the combined subjective \GTenToneHL{probability}{probability} of strong or moderate
	market conditions. The \GTenToneHL{probability}{probability} of strong market conditions is $\frac{\GTenToneHL{probability}{p}}{2}$,
	the \GTenToneHL{probability}{probability} of moderate market conditions is $\frac{\GTenToneHL{probability}{p}}{2}$, and the
	\GTenToneHL{probability}{probability} of weak market conditions is $\GTenToneHL{complement}{1-p}$. All \GTenToneHL{payoff}{payoffs} are expressed in
	thousands of dollars.
	
	Using the \GTenToneHL{expected}{Expected} \GTenToneHL{expected}{Value} criterion, determine which \GTenToneHL{alternative}{alternative} is preferred
	for different values of the subjective \GTenToneHL{probability}{probability} $\GTenToneHL{probability}{p}$. Identify the threshold
	\GTenToneHL{probability}{probability} at which the manager is indifferent between the two \GTenToneHL{alternative}{alternatives}.
	
	\vspace{-4mm}
	
	\begin{center}
		\renewcommand{\arraystretch}{1.3}
		\begin{tabular}{lccc}
			\toprule
			Alternative
			& \begin{tabular}{c}Strong market $\GTenToneHL{state}{S_1}$ \\ $\GTenToneHL{probability}{P}(\GTenToneHL{state}{S_1})=\frac{\GTenToneHL{probability}{p}}{2}$\end{tabular}
			& \begin{tabular}{c}Moderate market $\GTenToneHL{state}{S_2}$ \\ $\GTenToneHL{probability}{P}(\GTenToneHL{state}{S_2})=\frac{\GTenToneHL{probability}{p}}{2}$\end{tabular}
			& \begin{tabular}{c}Weak market $\GTenToneHL{state}{S_3}$ \\ $\GTenToneHL{probability}{P}(\GTenToneHL{state}{S_3})=\GTenToneHL{complement}{1-p}$\end{tabular}
			\\
			\midrule
			Automate production ($\GTenToneHL{alternative}{A_1}$) & 90 & 70 & 20 \\
			Keep current system ($\GTenToneHL{alternative}{A_2}$) & 55 & 45 & 40 \\
			\bottomrule
		\end{tabular}
	\end{center}
	
	\vspace{-5mm}
	
	% ============================================================
	% SOLUTION 2x3
	% ============================================================
	
	\subsection{Solution}
	
	\textbf{Step 1: Identify the \GTenToneHL{alternative}{alternatives}, \GTenToneHL{state}{states} \GTenToneHL{state}{of} \GTenToneHL{state}{nature}, and \GTenToneHL{probability}{probabilities}}
	
	Alternatives:
	$\GTenToneHL{alternative}{A_1}$ = Automate production,
	$\GTenToneHL{alternative}{A_2}$ = Keep current system
	
	States of nature:
	$\GTenToneHL{state}{S_1}$ = Strong market conditions,
	$\GTenToneHL{state}{S_2}$ = Moderate market conditions,
	$\GTenToneHL{state}{S_3}$ = Weak market conditions
	
	Probabilities:
	\[
	\GTenToneHL{probability}{P}(\GTenToneHL{state}{S_1})=\frac{\GTenToneHL{probability}{p}}{2},
	\qquad
	\GTenToneHL{probability}{P}(\GTenToneHL{state}{S_2})=\frac{\GTenToneHL{probability}{p}}{2},
	\qquad
	\GTenToneHL{probability}{P}(\GTenToneHL{state}{S_3})=\GTenToneHL{complement}{1-p}.
	\]
	
	Since
	\[
	\frac{\GTenToneHL{probability}{p}}{2}+\frac{\GTenToneHL{probability}{p}}{2}+(\GTenToneHL{complement}{1-p})=1,
	\]
	and all \GTenToneHL{probability}{probabilities} must be nonnegative, the feasible interval is
	\[
	0\leq \GTenToneHL{probability}{p}\leq1.
	\]
	
	\vspace{-3mm}
	
	\noindent
	\setlength{\columnseprule}{0.4pt}
	\begin{multicols}{2}
		\vspace{0pt}
		
		\textbf{Step 2: Calculate the \GTenToneHL{expected}{expected} \GTenToneHL{expected}{value} of $\GTenToneHL{alternative}{A_1}$}
		
		For the Automate production \GTenToneHL{alternative}{alternative}:
		\[
		\GTenToneHL{expected}{EV}(\GTenToneHL{alternative}{A_1})
		=
		\frac{\GTenToneHL{probability}{p}}{2}(90)
		+
		\frac{\GTenToneHL{probability}{p}}{2}(70)
		+
		(\GTenToneHL{complement}{1-p})(20)
		\]
		\[
		=
		45\GTenToneHL{probability}{p}+35\GTenToneHL{probability}{p}+20-20\GTenToneHL{probability}{p}
		=
		\GTenToneHL{probability}{60}\GTenToneHL{probability}{p}+20.
		\]
		
		Thus, the expected \GTenToneHL{payoff}{payoff} from automating production is
		$60\GTenToneHL{probability}{p}+20$ thousand dollars.
		
		\columnbreak
		
		\textbf{Step 3: Calculate the \GTenToneHL{expected}{expected} \GTenToneHL{expected}{value} of $\GTenToneHL{alternative}{A_2}$}
		
		For the Keep current system \GTenToneHL{alternative}{alternative}:
		\[
		\GTenToneHL{expected}{EV}(\GTenToneHL{alternative}{A_2})
		=
		\frac{\GTenToneHL{probability}{p}}{2}(55)
		+
		\frac{\GTenToneHL{probability}{p}}{2}(45)
		+
		(\GTenToneHL{complement}{1-p})(40)
		\]
		\[
		=
		27.5\GTenToneHL{probability}{p}+22.5\GTenToneHL{probability}{p}+40-40\GTenToneHL{probability}{p}
		=
		\GTenToneHL{probability}{10}\GTenToneHL{probability}{p}+40.
		\]
		
		Thus, the expected \GTenToneHL{payoff}{payoff} from keeping the current system is
		$10\GTenToneHL{probability}{p}+40$ thousand dollars.
		
	\end{multicols}
	
	\EMVGraph{
		{60}{20}{blue}{$\GTenToneHL{expected}{EMV}(\GTenToneHL{alternative}{A_1})=\GTenToneHL{expected}{60}\GTenToneHL{probability}{p}+20$},
		{10}{40}{red}{$\GTenToneHL{expected}{EMV}(\GTenToneHL{alternative}{A_2})=\GTenToneHL{expected}{10}\GTenToneHL{probability}{p}+40$}
	}
	
	\setlength{\columnseprule}{0pt}
	
	\textbf{Step 4: Find the threshold \GTenToneHL{probability}{probability}}
	
	The threshold occurs when the two \GTenToneHL{alternative}{alternatives} have the same \GTenToneHL{expected}{expected} \GTenToneHL{expected}{value}:
	\[
	\GTenToneHL{expected}{EV}(\GTenToneHL{alternative}{A_1})=\GTenToneHL{expected}{EV}(\GTenToneHL{alternative}{A_2})
	\qquad\longrightarrow\qquad
	60\GTenToneHL{probability}{p}+20=10\GTenToneHL{probability}{p}+40
	\]
	\[
	\qquad\longrightarrow\qquad
	50\GTenToneHL{probability}{p}=\GTenToneHL{probability}{20}
	\qquad\longrightarrow\qquad
	\GTenToneHL{probability}{p}=\frac{2}{5}=0.4.
	\]
	
	The threshold lies within the feasible interval $0\leq \GTenToneHL{probability}{p}\leq1$.
	
	\textbf{Step 5: Perform the sensitivity analysis}
	
	Compare the expected values:
	\[
	\GTenToneHL{expected}{EV}(\GTenToneHL{alternative}{A_1})-\GTenToneHL{expected}{EV}(\GTenToneHL{alternative}{A_2})
	=
	(60\GTenToneHL{probability}{p}+20)-(10\GTenToneHL{probability}{p}+40)
	=
	50\GTenToneHL{probability}{p}-20.
	\]
	
	Therefore,
	\[
	\boxed{
		\begin{aligned}
			\GTenToneHL{alternative}{A_2} &\text{ is preferred if } 0\leq \GTenToneHL{probability}{p}<0.4,\\
			\GTenToneHL{alternative}{A_1} &\text{ and }\GTenToneHL{alternative}{A_2}\text{ are equivalent if }\GTenToneHL{probability}{p}=\GTenToneHL{probability}{0.4},\\
			\GTenToneHL{alternative}{A_1} &\text{ is preferred if }0.4<\GTenToneHL{probability}{p}\leq1.
		\end{aligned}
	}
	\]
	
	The sensitivity analysis shows that keeping the current production system is
	preferred for lower values of $\GTenToneHL{probability}{p}$. Once $\GTenToneHL{probability}{p}$ exceeds $0.4$, investing in the
	automated production system has the higher expected \GTenToneHL{payoff}{payoff}.
	
	\textbf{Extended \GTenToneHL{payoff}{payoff} table}
	
	\begin{center}
		\renewcommand{\arraystretch}{1.5}
		\setlength{\tabcolsep}{4pt}
		
		\begin{tabular}{lcccc}
			\toprule
			Alternative
			& \begin{tabular}{c}Strong market $\GTenToneHL{state}{S_1}$ \\ $\GTenToneHL{probability}{P}(\GTenToneHL{state}{S_1})=\frac{\GTenToneHL{probability}{p}}{2}$\end{tabular}
			& \begin{tabular}{c}Moderate market $\GTenToneHL{state}{S_2}$ \\ $\GTenToneHL{probability}{P}(\GTenToneHL{state}{S_2})=\frac{\GTenToneHL{probability}{p}}{2}$\end{tabular}
			& \begin{tabular}{c}Weak market $\GTenToneHL{state}{S_3}$ \\ $\GTenToneHL{probability}{P}(\GTenToneHL{state}{S_3})=\GTenToneHL{complement}{1-p}$\end{tabular}
			& \GTenToneHL{expected}{Expected} \GTenToneHL{expected}{Value}
			\\
			\midrule
			
			$\GTenToneHL{alternative}{A_1}$: Automate production
			& $90$ & $70$ & $20$
			& $\frac{\GTenToneHL{probability}{p}}{2}(90)+\frac{\GTenToneHL{probability}{p}}{2}(70)+(\GTenToneHL{complement}{1-p})(20)=60\GTenToneHL{probability}{p}+20$
			\\
			
			$\GTenToneHL{alternative}{A_2}$: Keep current system
			& $55$ & $45$ & $40$
			& $\frac{\GTenToneHL{probability}{p}}{2}(55)+\frac{\GTenToneHL{probability}{p}}{2}(45)+(\GTenToneHL{complement}{1-p})(40)=10\GTenToneHL{probability}{p}+40$
			\\
			
			\bottomrule
		\end{tabular}
	\end{center}
	
	\section{3x2 - Technology company}
	
	A technology company is deciding which product strategy to adopt for the launch
	of a new device. The manager is uncertain about the level of demand during the
	launch period. Let $\GTenToneHL{probability}{p}$ represent the subjective \GTenToneHL{probability}{probability} of high demand.
	Therefore, the \GTenToneHL{probability}{probability} of low demand is $\GTenToneHL{complement}{1-p}$. All \GTenToneHL{payoff}{payoffs} are expressed
	in thousands of dollars.
	
	Using the \GTenToneHL{expected}{Expected} \GTenToneHL{expected}{Value} criterion, determine how the preferred decision
	changes as the subjective \GTenToneHL{probability}{probability} $\GTenToneHL{probability}{p}$ changes. Identify the threshold
	\GTenToneHL{probability}{probabilities} at which the manager is indifferent between the \GTenToneHL{alternative}{alternatives}.
	
	\vspace{-2mm}
	
	\begin{center}
		\renewcommand{\arraystretch}{1.3}
		\begin{tabular}{lcc}
			\toprule
			Alternative
			& \begin{tabular}{c}
				High demand $\GTenToneHL{state}{S_1}$ \\
				$\GTenToneHL{probability}{P}(\GTenToneHL{state}{S_1})=\GTenToneHL{probability}{p}$
			\end{tabular}
			& \begin{tabular}{c}
				Low demand $\GTenToneHL{state}{S_2}$ \\
				$\GTenToneHL{probability}{P}(\GTenToneHL{state}{S_2})=\GTenToneHL{complement}{1-p}$
			\end{tabular}
			\\
			\midrule
			Premium version ($\GTenToneHL{alternative}{A_1}$) & 80 & -20 \\
			Standard version ($\GTenToneHL{alternative}{A_2}$) & 60 & 18 \\
			Budget version ($\GTenToneHL{alternative}{A_3}$) & 40 & 25 \\
			\bottomrule
		\end{tabular}
	\end{center}
	
	\vspace{-5mm}
	
	\subsection{Solution}
	
	\textbf{Step 1: Identify the \GTenToneHL{alternative}{alternatives}, \GTenToneHL{state}{states} \GTenToneHL{state}{of} \GTenToneHL{state}{nature}, and \GTenToneHL{probability}{probabilities}}
	
	Alternatives:
	$\GTenToneHL{alternative}{A_1}$ = Premium version,
	$\GTenToneHL{alternative}{A_2}$ = Standard version,
	$\GTenToneHL{alternative}{A_3}$ = Budget version
	
	States of nature:
	$\GTenToneHL{state}{S_1}$ = High demand,
	$\GTenToneHL{state}{S_2}$ = Low demand
	
	Probabilities:
	$
	\GTenToneHL{probability}{P}(\GTenToneHL{state}{S_1})=\GTenToneHL{probability}{p},
	\GTenToneHL{probability}{P}(\GTenToneHL{state}{S_2})=\GTenToneHL{complement}{1-p}
	$
	with
	$
	0\leq \GTenToneHL{probability}{p}\leq1.
	$
	
	\vspace{-3mm}
	
	\noindent
	\setlength{\columnseprule}{0.4pt}
	\begin{multicols}{2}
		\vspace{0pt}
		
		\textbf{Step 2: Calculate the \GTenToneHL{expected}{expected} \GTenToneHL{expected}{value} of $\GTenToneHL{alternative}{A_1}$}
		
		For the Premium version \GTenToneHL{alternative}{alternative}:
		\[
		\GTenToneHL{expected}{EV}(\GTenToneHL{alternative}{A_1})
		=
		\GTenToneHL{probability}{p}(80)
		+
		(\GTenToneHL{complement}{1-p})(-20)
		\]
		\[
		=
		80\GTenToneHL{probability}{p}-20+20\GTenToneHL{probability}{p}
		\]
		\[
		=
		100\GTenToneHL{probability}{p}-20.
		\]
		
		Thus, the expected \GTenToneHL{payoff}{payoff} from the Premium version is
		$100\GTenToneHL{probability}{p}-20$ thousand dollars.
		
		\textbf{Step 3: Calculate the \GTenToneHL{expected}{expected} \GTenToneHL{expected}{value} of $\GTenToneHL{alternative}{A_2}$}
		
		For the Standard version \GTenToneHL{alternative}{alternative}:
		\[
		\GTenToneHL{expected}{EV}(\GTenToneHL{alternative}{A_2})
		=
		\GTenToneHL{probability}{p}(60)
		+
		(\GTenToneHL{complement}{1-p})(18)
		\]
		\[
		=
		60\GTenToneHL{probability}{p}+18-18\GTenToneHL{probability}{p}
		\]
		\[
		=
		42\GTenToneHL{probability}{p}+18.
		\]
		
		Thus, the expected \GTenToneHL{payoff}{payoff} from the Standard version is
		$42\GTenToneHL{probability}{p}+18$ thousand dollars.
		
		\columnbreak
		
		\textbf{Step 4: Calculate the \GTenToneHL{expected}{expected} \GTenToneHL{expected}{value} of $\GTenToneHL{alternative}{A_3}$}
		
		For the Budget version \GTenToneHL{alternative}{alternative}:
		\[
		\GTenToneHL{expected}{EV}(\GTenToneHL{alternative}{A_3})
		=
		\GTenToneHL{probability}{p}(40)
		+
		(\GTenToneHL{complement}{1-p})(25)
		\]
		\[
		=
		40\GTenToneHL{probability}{p}+25-25\GTenToneHL{probability}{p}
		\]
		\[
		=
		15\GTenToneHL{probability}{p}+25.
		\]
		
		Thus, the expected \GTenToneHL{payoff}{payoff} from the Budget version is
		$15\GTenToneHL{probability}{p}+25$ thousand dollars.
	\end{multicols}
	
	The thresholds occur when two \GTenToneHL{alternative}{alternatives} have the same \GTenToneHL{expected}{expected} \GTenToneHL{expected}{value}.
	
	\noindent
	\setlength{\columnseprule}{0.4pt}
	\begin{multicols}{2}
		\textbf{Step 5: Find the threshold \GTenToneHL{probability}{probabilities}}
		
		For $\GTenToneHL{alternative}{A_2}$ and $\GTenToneHL{alternative}{A_3}$:
		\[
		\GTenToneHL{expected}{EV}(\GTenToneHL{alternative}{A_2})=\GTenToneHL{expected}{EV}(\GTenToneHL{alternative}{A_3})
		\]
		\[
		42\GTenToneHL{probability}{p}+18=15\GTenToneHL{probability}{p}+25
		\]
		\[
		27\GTenToneHL{probability}{p}=\GTenToneHL{probability}{7}
		\]
		\[
		\GTenToneHL{probability}{p}=\frac{7}{27}
		\approx0.259.
		\]
		
		\columnbreak
		
		\textbf{Step 6: Find the second threshold \GTenToneHL{probability}{probability}}
		
		For $\GTenToneHL{alternative}{A_1}$ and $\GTenToneHL{alternative}{A_2}$:
		\[
		\GTenToneHL{expected}{EV}(\GTenToneHL{alternative}{A_1})=\GTenToneHL{expected}{EV}(\GTenToneHL{alternative}{A_2})
		\]
		\[
		100\GTenToneHL{probability}{p}-20=42\GTenToneHL{probability}{p}+18
		\]
		\[
		58\GTenToneHL{probability}{p}=\GTenToneHL{probability}{38}
		\]
		\[
		\GTenToneHL{probability}{p}=\frac{38}{58}
		=\frac{19}{29}
		\approx0.655.
		\]
		
	\end{multicols}
	
	Thus, the manager is indifferent between $\GTenToneHL{alternative}{A_2}$ and $\GTenToneHL{alternative}{A_3}$ when the
	subjective \GTenToneHL{probability}{probability} of high demand is approximately $25.9\%$,
	and indifferent between $\GTenToneHL{alternative}{A_1}$ and $\GTenToneHL{alternative}{A_2}$ when it is approximately
	$65.5\%$.
	
	\textbf{Step 7: Perform the sensitivity analysis}
	
	Compare the expected values on each side of the thresholds.
	
	$$
	\begin{aligned}
		\GTenToneHL{probability}{p}<0.259
		&\quad\Rightarrow\quad
		\GTenToneHL{expected}{EV}(\GTenToneHL{alternative}{A_3})>\GTenToneHL{expected}{EV}(\GTenToneHL{alternative}{A_2}),\GTenToneHL{expected}{EV}(\GTenToneHL{alternative}{A_1})
		\quad\Rightarrow\quad
		\text{choose }\GTenToneHL{alternative}{A_3},
		\\
		\GTenToneHL{probability}{p}=\GTenToneHL{probability}{0.259}
		&\quad\Rightarrow\quad
		\GTenToneHL{expected}{EV}(\GTenToneHL{alternative}{A_2})=\GTenToneHL{expected}{EV}(\GTenToneHL{alternative}{A_3})
		\quad\Rightarrow\quad
		\text{indifferent between }\GTenToneHL{alternative}{A_2}\text{ and }\GTenToneHL{alternative}{A_3},
		\\
		0.259<\GTenToneHL{probability}{p}<0.655
		&\quad\Rightarrow\quad
		\GTenToneHL{expected}{EV}(\GTenToneHL{alternative}{A_2})>\GTenToneHL{expected}{EV}(\GTenToneHL{alternative}{A_1}),\GTenToneHL{expected}{EV}(\GTenToneHL{alternative}{A_3})
		\quad\Rightarrow\quad
		\text{choose }\GTenToneHL{alternative}{A_2},
		\\
		\GTenToneHL{probability}{p}=\GTenToneHL{probability}{0.655}
		&\quad\Rightarrow\quad
		\GTenToneHL{expected}{EV}(\GTenToneHL{alternative}{A_1})=\GTenToneHL{expected}{EV}(\GTenToneHL{alternative}{A_2})
		\quad\Rightarrow\quad
		\text{indifferent between }\GTenToneHL{alternative}{A_1}\text{ and }\GTenToneHL{alternative}{A_2},
		\\
		\GTenToneHL{probability}{p}>0.655
		&\quad\Rightarrow\quad
		\GTenToneHL{expected}{EV}(\GTenToneHL{alternative}{A_1})>\GTenToneHL{expected}{EV}(\GTenToneHL{alternative}{A_2}),\GTenToneHL{expected}{EV}(\GTenToneHL{alternative}{A_3})
		\quad\Rightarrow\quad
		\text{choose }\GTenToneHL{alternative}{A_1}.
	\end{aligned}
	$$
	
	Therefore, the preferred decision depends on the subjective \GTenToneHL{probability}{probability}
	of high demand:
	
	$$
	\boxed{
		\begin{cases}
			\GTenToneHL{alternative}{A_3}, & 0\leq \GTenToneHL{probability}{p}<0.259,\\
			\text{indifferent between }\GTenToneHL{alternative}{A_2}\text{ and }\GTenToneHL{alternative}{A_3},
			& \GTenToneHL{probability}{p}=\GTenToneHL{probability}{0.259},\\
			\GTenToneHL{alternative}{A_2}, & 0.259<\GTenToneHL{probability}{p}<0.655,\\
			\text{indifferent between }\GTenToneHL{alternative}{A_1}\text{ and }\GTenToneHL{alternative}{A_2},
			& \GTenToneHL{probability}{p}=\GTenToneHL{probability}{0.655},\\
			\GTenToneHL{alternative}{A_1}, & 0.655<\GTenToneHL{probability}{p}\leq1.
	\end{cases}}
	$$
	
	The sensitivity analysis shows that the company should choose the Budget
	version when the subjective \GTenToneHL{probability}{probability} of high demand is below approximately
	$25.9\%$. Between $25.9\%$ and $65.5\%$, the Standard version has the highest
	\GTenToneHL{expected}{expected} \GTenToneHL{expected}{value}. When the \GTenToneHL{probability}{probability} of high demand exceeds $65.5\%$, the
	Premium version becomes the preferred decision.
	
	\textbf{Extended \GTenToneHL{payoff}{payoff} table}
	
	\begin{center}
		\renewcommand{\arraystretch}{1.5}
		\setlength{\tabcolsep}{3pt}
		
		\begin{tabular}{lccc}
			\toprule
			Alternative
			& \begin{tabular}{c}
				High demand $\GTenToneHL{state}{S_1}$ \\
				$\GTenToneHL{probability}{P}(\GTenToneHL{state}{S_1})=\GTenToneHL{probability}{p}$
			\end{tabular}
			& \begin{tabular}{c}
				Low demand $\GTenToneHL{state}{S_2}$ \\
				$\GTenToneHL{probability}{P}(\GTenToneHL{state}{S_2})=\GTenToneHL{complement}{1-p}$
			\end{tabular}
			& \GTenToneHL{expected}{Expected} \GTenToneHL{expected}{Value}
			\\
			\midrule
			
			$\GTenToneHL{alternative}{A_1}$: Premium version
			& $80$
			& $-20$
			& \begin{tabular}{c}
				$\GTenToneHL{probability}{p}(80)+(\GTenToneHL{complement}{1-p})(-20)$ \\
				$=100\GTenToneHL{probability}{p}-20$
			\end{tabular}
			\\
			$\GTenToneHL{alternative}{A_2}$: Standard version
			& $60$
			& $18$
			& \begin{tabular}{c}
				$\GTenToneHL{probability}{p}(60)+(\GTenToneHL{complement}{1-p})(18)$ \\
				$=42\GTenToneHL{probability}{p}+18$
			\end{tabular}
			\\
			$\GTenToneHL{alternative}{A_3}$: Budget version
			& $40$
			& $25$
			& \begin{tabular}{c}
				$\GTenToneHL{probability}{p}(40)+(\GTenToneHL{complement}{1-p})(25)$ \\
				$=15\GTenToneHL{probability}{p}+25$
			\end{tabular}
			\\
			\bottomrule
		\end{tabular}
	\end{center}
	
	\vspace{2mm}
	
	\textbf{\GTenToneHL{expected}{Expected} \GTenToneHL{expected}{Value} graph}
	
	\begin{center}
		\EMVGraph{
			{100}{-20}{blue}{$\GTenToneHL{expected}{EV}(\GTenToneHL{alternative}{A_1})=\GTenToneHL{expected}{100}\GTenToneHL{probability}{p}-20$},
			{42}{18}{red}{$\GTenToneHL{expected}{EV}(\GTenToneHL{alternative}{A_2})=\GTenToneHL{expected}{42}\GTenToneHL{probability}{p}+18$},
			{15}{25}{green}{$\GTenToneHL{expected}{EV}(\GTenToneHL{alternative}{A_3})=\GTenToneHL{expected}{15}\GTenToneHL{probability}{p}+25$}
		}
		\end{center}
	
	\section{3x2 - Renewable energy company}

	A renewable energy company must choose one of three solar-farm designs.
	The return depends on whether electricity demand is high ($\GTenToneHL{state}{S_1}$) or low
	($\GTenToneHL{state}{S_2}$). Let $\GTenToneHL{probability}{P}(\GTenToneHL{state}{S_1})=\GTenToneHL{probability}{p}$ and $\GTenToneHL{probability}{P}(\GTenToneHL{state}{S_2})=\GTenToneHL{complement}{1-p}$. Payoffs are in thousands of
	dollars.

	\begin{center}
		\renewcommand{\arraystretch}{1.3}
		\begin{tabular}{lcc}
			\toprule
			Alternative & High demand $\GTenToneHL{state}{S_1}$ & Low demand $\GTenToneHL{state}{S_2}$ \\
			\midrule
			Large solar farm ($\GTenToneHL{alternative}{A_1}$) & \GTenToneHL{payoff}{90} & $-30$ \\
			Medium solar farm ($\GTenToneHL{alternative}{A_2}$) & \GTenToneHL{payoff}{60} & \GTenToneHL{payoff}{15} \\
			Small solar farm ($\GTenToneHL{alternative}{A_3}$) & \GTenToneHL{payoff}{35} & \GTenToneHL{payoff}{30} \\
			\bottomrule
		\end{tabular}
	\end{center}

	Using the \GTenToneHL{expected}{Expected} \GTenToneHL{expected}{Value} criterion, determine the preferred design for
	each value of $\GTenToneHL{probability}{p}$ and identify the threshold \GTenToneHL{probability}{probabilities}.

	\subsection{Solution}

	The expected-value functions are
	\[
		\GTenToneHL{expected}{EV}(\GTenToneHL{alternative}{A_1})=\GTenToneHL{expected}{120}\GTenToneHL{probability}{p}-30,\qquad
		\GTenToneHL{expected}{EV}(\GTenToneHL{alternative}{A_2})=\GTenToneHL{expected}{45}\GTenToneHL{probability}{p}+15,\qquad
		\GTenToneHL{expected}{EV}(\GTenToneHL{alternative}{A_3})=\GTenToneHL{expected}{5}\GTenToneHL{probability}{p}+30.
	\]
	The relevant intersections are
	\[
		\GTenToneHL{expected}{EV}(\GTenToneHL{alternative}{A_2})=\GTenToneHL{expected}{EV}(\GTenToneHL{alternative}{A_3})\Rightarrow \GTenToneHL{probability}{p}=\frac{3}{8},
		\qquad
		\GTenToneHL{expected}{EV}(\GTenToneHL{alternative}{A_1})=\GTenToneHL{expected}{EV}(\GTenToneHL{alternative}{A_2})\Rightarrow \GTenToneHL{probability}{p}=\frac{3}{5}.
	\]
	(The intersection of $\GTenToneHL{alternative}{A_1}$ and $\GTenToneHL{alternative}{A_3}$ is not a decision threshold because
	$\GTenToneHL{alternative}{A_2}$ has a greater \GTenToneHL{expected}{expected} \GTenToneHL{expected}{value} there.) Therefore,
	\[
		\boxed{
		\begin{cases}
			\GTenToneHL{alternative}{A_3}, & 0\leq \GTenToneHL{probability}{p}<\frac{3}{8},\\
			\text{indifferent between }\GTenToneHL{alternative}{A_2}\text{ and }\GTenToneHL{alternative}{A_3},
				& \GTenToneHL{probability}{p}=\frac{3}{8},\\
			\GTenToneHL{alternative}{A_2}, & \frac{3}{8}<\GTenToneHL{probability}{p}<\frac{3}{5},\\
			\text{indifferent between }\GTenToneHL{alternative}{A_1}\text{ and }\GTenToneHL{alternative}{A_2},
				& \GTenToneHL{probability}{p}=\frac{3}{5},\\
			\GTenToneHL{alternative}{A_1}, & \frac{3}{5}<\GTenToneHL{probability}{p}\leq1.
		\end{cases}}
	\]

	\textbf{Extended \GTenToneHL{payoff}{payoff} table}
	\begin{center}
		\renewcommand{\arraystretch}{1.4}
		\begin{tabular}{lccc}
			\toprule
			Alternative & $\GTenToneHL{state}{S_1}$ $(\GTenToneHL{probability}{p})$ & $\GTenToneHL{state}{S_2}$ $(\GTenToneHL{complement}{1-p})$ & \GTenToneHL{expected}{Expected} \GTenToneHL{expected}{Value} \\
			\midrule
			$\GTenToneHL{alternative}{A_1}$: Large & \GTenToneHL{probability}{90} & $-30$ & $120\GTenToneHL{probability}{p}-30$ \\
			$\GTenToneHL{alternative}{A_2}$: Medium & \GTenToneHL{probability}{60} & \GTenToneHL{payoff}{15} & $45\GTenToneHL{probability}{p}+15$ \\
			$\GTenToneHL{alternative}{A_3}$: Small & \GTenToneHL{probability}{35} & \GTenToneHL{payoff}{30} & $5\GTenToneHL{probability}{p}+30$ \\
			\bottomrule
		\end{tabular}
	\end{center}

	\textbf{\GTenToneHL{expected}{Expected} \GTenToneHL{expected}{Value} graph}
	\begin{center}
		\EMVGraph{
			{120}{-30}{blue}{$\GTenToneHL{expected}{EV}(\GTenToneHL{alternative}{A_1})=\GTenToneHL{expected}{120}\GTenToneHL{probability}{p}-30$},
			{45}{15}{red}{$\GTenToneHL{expected}{EV}(\GTenToneHL{alternative}{A_2})=\GTenToneHL{expected}{45}\GTenToneHL{probability}{p}+15$},
			{5}{30}{green}{$\GTenToneHL{expected}{EV}(\GTenToneHL{alternative}{A_3})=\GTenToneHL{expected}{5}\GTenToneHL{probability}{p}+30$}
		}
	\end{center}

	\section{3x2 - Restaurant chain}

	A restaurant chain is choosing a format for a new location. Customer
	traffic may be high ($\GTenToneHL{state}{S_1}$) or low ($\GTenToneHL{state}{S_2}$), with \GTenToneHL{probability}{probabilities} $\GTenToneHL{probability}{p}$ and
	$\GTenToneHL{complement}{1-p}$, respectively. Payoffs are in thousands of dollars.

	\begin{center}
		\renewcommand{\arraystretch}{1.3}
		\begin{tabular}{lcc}
			\toprule
			Alternative & High traffic $\GTenToneHL{state}{S_1}$ & Low traffic $\GTenToneHL{state}{S_2}$ \\
			\midrule
			Flagship restaurant ($\GTenToneHL{alternative}{A_1}$) & \GTenToneHL{payoff}{110} & $-40$ \\
			Neighborhood restaurant ($\GTenToneHL{alternative}{A_2}$) & \GTenToneHL{payoff}{70} & \GTenToneHL{payoff}{20} \\
			Pop-up restaurant ($\GTenToneHL{alternative}{A_3}$) & \GTenToneHL{payoff}{45} & \GTenToneHL{payoff}{35} \\
			\bottomrule
		\end{tabular}
	\end{center}

	Using the \GTenToneHL{expected}{Expected} \GTenToneHL{expected}{Value} criterion, determine the preferred format for
	each value of $\GTenToneHL{probability}{p}$ and identify the threshold \GTenToneHL{probability}{probabilities}.

	\subsection{Solution}

	The expected-value functions are
	\[
		\GTenToneHL{expected}{EV}(\GTenToneHL{alternative}{A_1})=\GTenToneHL{expected}{150}\GTenToneHL{probability}{p}-40,\qquad
		\GTenToneHL{expected}{EV}(\GTenToneHL{alternative}{A_2})=\GTenToneHL{expected}{50}\GTenToneHL{probability}{p}+20,\qquad
		\GTenToneHL{expected}{EV}(\GTenToneHL{alternative}{A_3})=\GTenToneHL{expected}{10}\GTenToneHL{probability}{p}+35.
	\]
	The relevant intersections are
	\[
		\GTenToneHL{expected}{EV}(\GTenToneHL{alternative}{A_2})=\GTenToneHL{expected}{EV}(\GTenToneHL{alternative}{A_3})\Rightarrow \GTenToneHL{probability}{p}=\frac{3}{8},
		\qquad
		\GTenToneHL{expected}{EV}(\GTenToneHL{alternative}{A_1})=\GTenToneHL{expected}{EV}(\GTenToneHL{alternative}{A_2})\Rightarrow \GTenToneHL{probability}{p}=\frac{3}{5}.
	\]
	Although $\GTenToneHL{alternative}{A_1}$ and $\GTenToneHL{alternative}{A_3}$ also intersect, $\GTenToneHL{alternative}{A_2}$ has a greater expected
	value at that \GTenToneHL{probability}{probability}. The decision rule is
	\[
		\boxed{
		\begin{cases}
			\GTenToneHL{alternative}{A_3}, & 0\leq \GTenToneHL{probability}{p}<\frac{3}{8},\\
			\text{indifferent between }\GTenToneHL{alternative}{A_2}\text{ and }\GTenToneHL{alternative}{A_3},
				& \GTenToneHL{probability}{p}=\frac{3}{8},\\
			\GTenToneHL{alternative}{A_2}, & \frac{3}{8}<\GTenToneHL{probability}{p}<\frac{3}{5},\\
			\text{indifferent between }\GTenToneHL{alternative}{A_1}\text{ and }\GTenToneHL{alternative}{A_2},
				& \GTenToneHL{probability}{p}=\frac{3}{5},\\
			\GTenToneHL{alternative}{A_1}, & \frac{3}{5}<\GTenToneHL{probability}{p}\leq1.
		\end{cases}}
	\]

	\textbf{Extended \GTenToneHL{payoff}{payoff} table}
	\begin{center}
		\renewcommand{\arraystretch}{1.4}
		\begin{tabular}{lccc}
			\toprule
			Alternative & $\GTenToneHL{state}{S_1}$ $(\GTenToneHL{probability}{p})$ & $\GTenToneHL{state}{S_2}$ $(\GTenToneHL{complement}{1-p})$ & \GTenToneHL{expected}{Expected} \GTenToneHL{expected}{Value} \\
			\midrule
			$\GTenToneHL{alternative}{A_1}$: Flagship & \GTenToneHL{probability}{110} & $-40$ & $150\GTenToneHL{probability}{p}-40$ \\
			$\GTenToneHL{alternative}{A_2}$: Neighborhood & \GTenToneHL{probability}{70} & \GTenToneHL{payoff}{20} & $50\GTenToneHL{probability}{p}+20$ \\
			$\GTenToneHL{alternative}{A_3}$: Pop-up & \GTenToneHL{probability}{45} & \GTenToneHL{payoff}{35} & $10\GTenToneHL{probability}{p}+35$ \\
			\bottomrule
		\end{tabular}
	\end{center}

	\textbf{\GTenToneHL{expected}{Expected} \GTenToneHL{expected}{Value} graph}
	\begin{center}
		\EMVGraph{
			{150}{-40}{blue}{$\GTenToneHL{expected}{EV}(\GTenToneHL{alternative}{A_1})=\GTenToneHL{expected}{150}\GTenToneHL{probability}{p}-40$},
			{50}{20}{red}{$\GTenToneHL{expected}{EV}(\GTenToneHL{alternative}{A_2})=\GTenToneHL{expected}{50}\GTenToneHL{probability}{p}+20$},
			{10}{35}{green}{$\GTenToneHL{expected}{EV}(\GTenToneHL{alternative}{A_3})=\GTenToneHL{expected}{10}\GTenToneHL{probability}{p}+35$}
		}
	\end{center}

	\section{3x2 - Agricultural cooperative}

	An agricultural cooperative must choose a planting strategy before it
	knows whether the growing season will be favorable ($\GTenToneHL{state}{S_1}$) or unfavorable
	($\GTenToneHL{state}{S_2}$). Let $\GTenToneHL{probability}{P}(\GTenToneHL{state}{S_1})=\GTenToneHL{probability}{p}$ and $\GTenToneHL{probability}{P}(\GTenToneHL{state}{S_2})=\GTenToneHL{complement}{1-p}$. Payoffs are in thousands of
	dollars.

	\begin{center}
		\renewcommand{\arraystretch}{1.3}
		\begin{tabular}{lcc}
			\toprule
			Alternative & Favorable $\GTenToneHL{state}{S_1}$ & Unfavorable $\GTenToneHL{state}{S_2}$ \\
			\midrule
			High-intensity plan ($\GTenToneHL{alternative}{A_1}$) & \GTenToneHL{payoff}{95} & $-25$ \\
			Standard plan ($\GTenToneHL{alternative}{A_2}$) & \GTenToneHL{payoff}{65} & \GTenToneHL{payoff}{15} \\
			Resilient plan ($\GTenToneHL{alternative}{A_3}$) & \GTenToneHL{payoff}{40} & \GTenToneHL{payoff}{30} \\
			\bottomrule
		\end{tabular}
	\end{center}

	Using the \GTenToneHL{expected}{Expected} \GTenToneHL{expected}{Value} criterion, determine the preferred planting plan
	for each value of $\GTenToneHL{probability}{p}$ and identify the threshold \GTenToneHL{probability}{probabilities}.

	\subsection{Solution}

	The expected-value functions are
	\[
		\GTenToneHL{expected}{EV}(\GTenToneHL{alternative}{A_1})=\GTenToneHL{expected}{120}\GTenToneHL{probability}{p}-25,\qquad
		\GTenToneHL{expected}{EV}(\GTenToneHL{alternative}{A_2})=\GTenToneHL{expected}{50}\GTenToneHL{probability}{p}+15,\qquad
		\GTenToneHL{expected}{EV}(\GTenToneHL{alternative}{A_3})=\GTenToneHL{expected}{10}\GTenToneHL{probability}{p}+30.
	\]
	The relevant intersections are
	\[
		\GTenToneHL{expected}{EV}(\GTenToneHL{alternative}{A_2})=\GTenToneHL{expected}{EV}(\GTenToneHL{alternative}{A_3})\Rightarrow \GTenToneHL{probability}{p}=\frac{3}{8},
		\qquad
		\GTenToneHL{expected}{EV}(\GTenToneHL{alternative}{A_1})=\GTenToneHL{expected}{EV}(\GTenToneHL{alternative}{A_2})\Rightarrow \GTenToneHL{probability}{p}=\frac{4}{7}.
	\]
	The intersection of $\GTenToneHL{alternative}{A_1}$ and $\GTenToneHL{alternative}{A_3}$ is not a decision threshold because
	$\GTenToneHL{alternative}{A_2}$ has a greater \GTenToneHL{expected}{expected} \GTenToneHL{expected}{value} there. Thus,
	\[
		\boxed{
		\begin{cases}
			\GTenToneHL{alternative}{A_3}, & 0\leq \GTenToneHL{probability}{p}<\frac{3}{8},\\
			\text{indifferent between }\GTenToneHL{alternative}{A_2}\text{ and }\GTenToneHL{alternative}{A_3},
				& \GTenToneHL{probability}{p}=\frac{3}{8},\\
			\GTenToneHL{alternative}{A_2}, & \frac{3}{8}<\GTenToneHL{probability}{p}<\frac{4}{7},\\
			\text{indifferent between }\GTenToneHL{alternative}{A_1}\text{ and }\GTenToneHL{alternative}{A_2},
				& \GTenToneHL{probability}{p}=\frac{4}{7},\\
			\GTenToneHL{alternative}{A_1}, & \frac{4}{7}<\GTenToneHL{probability}{p}\leq1.
		\end{cases}}
	\]

	\textbf{Extended \GTenToneHL{payoff}{payoff} table}
	\begin{center}
		\renewcommand{\arraystretch}{1.4}
		\begin{tabular}{lccc}
			\toprule
			Alternative & $\GTenToneHL{state}{S_1}$ $(\GTenToneHL{probability}{p})$ & $\GTenToneHL{state}{S_2}$ $(\GTenToneHL{complement}{1-p})$ & \GTenToneHL{expected}{Expected} \GTenToneHL{expected}{Value} \\
			\midrule
			$\GTenToneHL{alternative}{A_1}$: High-intensity & \GTenToneHL{probability}{95} & $-25$ & $120\GTenToneHL{probability}{p}-25$ \\
			$\GTenToneHL{alternative}{A_2}$: Standard & \GTenToneHL{probability}{65} & \GTenToneHL{payoff}{15} & $50\GTenToneHL{probability}{p}+15$ \\
			$\GTenToneHL{alternative}{A_3}$: Resilient & \GTenToneHL{probability}{40} & \GTenToneHL{payoff}{30} & $10\GTenToneHL{probability}{p}+30$ \\
			\bottomrule
		\end{tabular}
	\end{center}

	\textbf{\GTenToneHL{expected}{Expected} \GTenToneHL{expected}{Value} graph}
	\begin{center}
		\EMVGraph{
			{120}{-25}{blue}{$\GTenToneHL{expected}{EV}(\GTenToneHL{alternative}{A_1})=\GTenToneHL{expected}{120}\GTenToneHL{probability}{p}-25$},
			{50}{15}{red}{$\GTenToneHL{expected}{EV}(\GTenToneHL{alternative}{A_2})=\GTenToneHL{expected}{50}\GTenToneHL{probability}{p}+15$},
			{10}{30}{green}{$\GTenToneHL{expected}{EV}(\GTenToneHL{alternative}{A_3})=\GTenToneHL{expected}{10}\GTenToneHL{probability}{p}+30$}
		}
	\end{center}

	\end{document}
